\subsection{Preliminaries on automorphisms of rings of formal power series}

LEMMA 2. Let \(f(\mathbf{X}_1, \mathbf{X}_2, \ldots, \mathbf{X}_n)\) be a formal power series \(\neq 0\) with coefficients in a ring E. There exist integers \(u(i) \geqslant 1\) ( \(1 \leqslant i \leqslant n - 1\) ) such that

\[
f (\mathrm{T} ^ {u (1)}, \dots , \mathrm{T} ^ {u (n - 1)}, \mathrm{T}) \neq 0.
\]

Suppose that integers \(u(i) \geqslant 1\) ( \(1 \leqslant i \leqslant k - 1\) ) are determined such that \(f(\mathbf{X}_n^{u(1)}, \ldots, \mathbf{X}_n^{u(k-1)}, \mathbf{X}_k, \ldots, \mathbf{X}_n) \neq 0\) . We shall determine an integer \(u(k) \geqslant 1\) such that

\[
f \left(\mathrm{X} _ {n} ^ {u (1)}, \dots , \mathrm{X} _ {n} ^ {u (k - 1)}, \mathrm{X} _ {n} ^ {u (k)}, \mathrm{X} _ {k + 1}, \dots , \mathrm{X} _ {n}\right) \neq 0.
\]

The lemma will then be proved by induction.

Observe that the series \(f(\mathbf{X}_n^{u(1)}, \ldots, \mathbf{X}_n^{u(k-1)}, \mathbf{X}_k, \ldots, \mathbf{X}_n)\) can be considered as a series in \(\mathbf{X}_k\) and \(\mathbf{X}_n\) with coefficients in \(\mathbf{E}[[\mathbf{X}_{k+1}, \ldots, \mathbf{X}_{n-1}]]\) . Thus we see that it suffices to establish the lemma for \(n = 2\) .

Therefore let

\[
f = \sum_ {i, j} e _ {i j} \mathbf {X} ^ {i} \mathbf {Y} ^ {j} \in \mathrm{E} [ [ \mathbf {X}, \mathbf {Y} ] ]
\]

where \(f \neq 0\) . Let \(G \subset N \times N\) be the non-empty set of ordered pairs \((i, j)\) such that \(e_{ij} \neq 0\) . Let \(N \times N\) be given the lexicographical ordering. Let \((c, d)\) be the least element of \(G\) . Choose an integer \(p > d\) . In the expansion of

\[
f (\mathbf {T} ^ {p}, \mathbf {T}) = \sum_ {(i, j) \in \mathbb {G}} e _ {i j} \mathbf {T} ^ {i p + j}
\]

we look for the terms of degree \(cp + d\) . If \(ip + j = cp + d\) , \(i \geqslant c + 1\) is impossible, for this would give

\[
i p + j \geqslant (c + 1) p + j \geqslant (c + 1) p > c p + d;
\]

nor is \(i < c\) possible, for \((c, d)\) is the least element of \(G\) ; therefore \(i = c\) and then \(j = d\) . The term of degree \(cp + d\) in \(f(T^p, T)\) is therefore \(e_{cd}T^{cp + d}\) . Since \(e_{cd} \neq 0, f(T^p, T) \neq 0\) . Whence the lemma.

In the ring \(\operatorname{E}[[\mathbf{X}_1, \ldots, \mathbf{X}_n]]\) , let \(a\) be the ideal of formal power series without constant term. If \(w_1, \ldots, w_n\) are elements of \(a\) , recall that the mapping \(f(\mathbf{X}_1, \ldots, \mathbf{X}_n) \mapsto f(w_1, \ldots, w_n)\) is the unique endomorphism \(s\) of the ring \(\operatorname{E}[[\mathbf{X}_1, \ldots, \mathbf{X}_n]]\) such that \(s(\mathbf{X}_i) = w_i\) for \(1 \leqslant i \leqslant n\) (Chapter III, § 4, no. 5, Proposition 6).

We take \(w_1 = \mathbf{X}_1 + \mathbf{X}_n^{u(1)}, \ldots, w_{n-1} = \mathbf{X}_{n-1} + \mathbf{X}_n^{u(n-1)}, w_n = \mathbf{X}_n\) , where the \(u(i)\) are integers \(\geqslant 1\) . Let \(s'\) be the endomorphism of \(\operatorname{E}[[\mathbf{X}_1, \ldots, \mathbf{X}_n]]\) which maps \(\mathbf{X}_1\) to \(\mathbf{X}_1 - \mathbf{X}_n^{u(1)}, \ldots, \mathbf{X}_{n-1}\) to \(\mathbf{X}_{n-1} - \mathbf{X}_n^{u(n-1)}\) and \(\mathbf{X}_n\) to \(\mathbf{X}_n\) . Then \(s'(s(\mathbf{X}_i)) = \mathbf{X}_i\) for \(1 \leqslant i \leqslant n\) and hence \(s' \circ s\) is the identity automorphism; similarly for \(s \circ s'\) . Hence \(s\) is an automorphism.

LEMMA 3. Let \(f\) be a non-zero element of \(\operatorname{E}[[\mathbf{X}_1, \ldots, \mathbf{X}_n]]\) . There exist integers \(u(i) \geqslant 1 (1 \leqslant i \leqslant n - 1)\) such that the automorphism \(s\) of \(\operatorname{E}\) defined by

\[
s (\mathbf {X} _ {i}) = \mathbf {X} _ {i} + \mathbf {X} _ {n} ^ {u (i)} \quad (1 \leqslant i \leqslant n - 1)
\]

and \(s(\mathbf{X}_n) = \mathbf{X}_n\) maps \(f\) to an element \(g\) such that \(g(0, \ldots, 0, \mathbf{X}_n) \neq 0\) .

\(g(0,\ldots,0,\mathbf{X}_n) = f(\mathbf{X}_n^{u(1)},\ldots ,\mathbf{X}_n^{u(n - 1)},\mathbf{X}_n)\) . Lemma 3 is therefore a consequence of Lemma 2.

\subsection{The preparation theorem}

In this no. A will denote a local ring, m its maximal ideal and k = A/m its residue field. Suppose that A is Hausdorff and complete with the m-adic topology. Let B = A{[}{[}X{]}{]}; it is local ring whose maximal ideal N is generated by m and X; with the \(\mathfrak{N}\) -adic topology, B is Hausdorff and complete (Chapter III, § 2, no. 6, Proposition 6).

For every formal power series

\[
f = \sum_ {i = 0} ^ {\infty} a _ {i} \mathbf {X} ^ {i} \in \mathbf {B},
\]

we write

\[
\bar {f} = \sum_ {i = 0} ^ {\infty} \bar {a} _ {i} X ^ {i} \in k [ [ X ] ],
\]

where \(\bar{a}_{i}\) denotes the canonical image of \(a_{i}\) in k. The series \(\bar{f}\) will be called the reduced series of f; if \(\bar{f} \neq 0\) , the order of \(\bar{f}\) (that is the least integer s such that \(a_{s} \notin m\) ) will be called the reduced order of f.

PROPOSITION 5. Let \(f \in \mathbf{B}\) be a series whose reduced series is non-zero. Let \(s\) denote its reduced order and \(\mathbf{M}\) the sub-A-module of \(\mathbf{B}\) with basis \(\{1, X, \ldots, X^{s-1}\}\) . Then \(\mathbf{B}\) is the direct sum of \(\mathbf{M}\) and \(f\mathbf{B}\) and \(f\) is not a divisor of zero in \(\mathbf{B}\) .

\begin{enumerate}
\def\labelenumi{(\alph{enumi})}
\tightlist
\item
  We show that \(f\mathbf{B} \cap \mathbf{M} = (0)\) . Suppose that there is a relation:
\end{enumerate}

\[
\left(\sum_ {i = 0} ^ {\infty} b _ {i} \mathrm{X} ^ {i}\right) f = r _ {0} + r _ {1} \mathrm{X} + \dots + r _ {s - 1} \mathrm{X} ^ {s - 1} \quad (b _ {i} \in \mathrm{A}, r _ {j} \in \mathrm{A}).\tag{3}
\]

We show that the \(b_{i}\) (and hence the \(r_{j}\) ) are all zero, which will prove in particular that \(f\) is not a divisor of zero in B. Since A is Hausdorff, it suffices to show that \(b_{i} \in \mathfrak{m}^{n}\) for all \(i \geqslant 0\) and all \(n \geqslant 0\) . It is obvious for \(n = 0\) . We shall argue by double induction: we shall assume that \(b_{i} \in \mathfrak{m}^{n-1}\) for all \(i\) and \(b_{i} \in \mathfrak{m}^{n}\) for \(i < k\) and show that this implies that \(b_{k} \in \mathfrak{m}^{n}\) . For this, we write \(f = \sum_{i=1}^{\infty} a_{i} X^{i}\) and compare the coefficients of \(\mathbf{X}^{s+k}\) in (3); then:

\[
\left(b _ {0} a _ {s + k} + \dots + b _ {k - 1} a _ {s + 1}\right) + b _ {k} a _ {s} + \left(b _ {k + 1} a _ {s - 1} + \dots + b _ {k + s} a _ {0}\right) = 0.\tag{4}
\]

The terms in the first bracket belong to \(\mathfrak{m}^n\) since the \(b_i \in \mathfrak{m}^n\) for \(i < k\) ; similarly for those in the second, since the \(b_i \in \mathfrak{m}^{n-1}\) for all \(i\) and the \(a_i \in \mathfrak{m}\) for \(i \leqslant s - 1\) . Hence \(b_k a_s \in \mathfrak{m}^n\) and, as \(a_s\) is an invertible element of A, \(b_k \in \mathfrak{m}^n\) .

\begin{enumerate}
\def\labelenumi{(\alph{enumi})}
\setcounter{enumi}{1}
\tightlist
\item
  We show that \(f\mathbf{B} + \mathbf{M} = \mathbf{B}\) . We write
\end{enumerate}

\[
g = a _ {s} + a _ {s + 1} \mathbf {X} + a _ {s + 2} \mathbf {X} ^ {2} + \dots ;
\]

it is an invertible element of B. Then

\[
f - \mathrm{X} ^ {s} g = a _ {0} + a _ {1} \mathrm{X} + \dots + a _ {s - 1} \mathrm{X} ^ {s - 1};
\]

if therefore we write \(fg^{-1} - \mathbf{X}^s = (f - \mathbf{X}^s g)g^{-1} = -h\) , the coefficients of \(h\) belong to \(m\) . Then let \(r\) be an element of B. By induction on \(n\) we define a sequence \((q^{(n)})\) of elements of B: we take \(q^{(0)}\) to be the unique series satisfying:

\[
r \equiv \mathbf {X} ^ {s} q ^ {(0)} \pmod {\mathbf {M}};\tag{5}
\]

writing \(h = \sum_{i=0}^{\infty} h_i X^i\) and \(q^{(n)} = \sum_{i=0}^{\infty} q_i^{(n)} X^i\) , the \(q_i^{(n)}\) are defined by:

\[
q _ {i} ^ {(n)} = \sum_ {j = 0} ^ {i + s} h _ {j} q _ {i + s - j} ^ {(n - 1)}.\tag{6}
\]

It follows immediately from (6) that:

\[
\mathbf {X} ^ {s} q ^ {(n)} \equiv h q ^ {(n - 1)} \pmod {\mathbf {M}}.\tag{7}
\]

As \(h_j \in \mathfrak{m}\) for all \(j\) , it also follows from (6), by induction on \(n\) , that \(q_i^{(n)} \in \mathfrak{m}^n\) for all \(i\) and all \(n\) . As \(A\) is complete, it follows that the series

\[
q ^ {(0)} + q ^ {(1)} + \dots + q ^ {(n)} + \dots
\]

converges to an element \(q\) of B. By (5) and (7),

\[
\mathbf {X} ^ {s} (q ^ {(0)} + q ^ {(1)} + \dots + q ^ {(n)}) \equiv r + h (q ^ {(0)} + \dots + q ^ {(n - 1)}) \pmod {\mathbf {M}}. \tag {8}
\]

As \(\mathbf{M}\) is closed, at the limit (8) gives \(r \equiv (\mathbf{X}^s - h)q\) (mod. \(\mathbf{M}\) ), that is

\[
r \in f g ^ {- 1} q + \mathbf {M} \subset f \mathbf {B} + \mathbf {M}.
\]

We may also use the results of Chapter III, §2 to show the relation \(B = fB + M\) (cf.~Exercise 12). The method followed here has the advantage of being applicable to convergent series.

COROLLARY. With the hypotheses and notation of Proposition 5, suppose that \(s \geqslant 1\) , so that \(f \in \mathbf{B}\mathfrak{m} + \mathbf{B}\mathbf{X}\) . Then the A-homomorphism \(h\) of \(\mathbf{B}' = \mathbf{A}[[\mathbf{T}]]\) to \(\mathbf{B} = \mathbf{A}[[\mathbf{X}]]\) such that \(h(\mathbf{T}) = f\) (Chapter III, §2, no. 9, Proposition 11 (a)) defines on \(\mathbf{B}\) a free \(\mathbf{B}'\) -module structure admitting \(\{1, \mathbf{X}, \ldots, \mathbf{X}^{s-1}\}\) as basis. In particular \(h\) is injective.

Let the B'-module B be given the (T)-adic filtration, which consists of the \(f^{n}B\) for \(n \geqslant 0\) (Chapter III, §2, no. 1). Then B/fB is a free module over the ring A = B'/TB' and the images of the \(X^{i}\) ( \(0 \leqslant i \leqslant s - 1\) ) in this A-module form a basis of it (Proposition 5); as moreover f is not a divisor of zero in B (Proposition 5), \(Bf^{n}/Bf^{n+1}\) is also a free (B'/TB')-module of rank s, so that condition (GR) of Chapter III, §2, no. 8 is satisfied (replacing A by B' and M by B). On the other hand, since B' is Hausdorff and complete with respect to the (T)-adic filtration and gr(B) is a finitely generated gr(B')-module by the above, it is seen first (Chapter III, §2, no. 9, Corollary 1 to Proposition 12) that B is a finitely generated B' module. The first assertion of the corollary then follows from Chapter III, §2, no. 9, Proposition 13. The second follows immediately from it.

DEFINITION 2. A polynomial \(\mathbf{F} \in \mathbf{A}[\mathbf{X}]\) is called distinguished if it is of the form

\[
\mathbf {F} = \mathbf {X} ^ {s} + a _ {s - 1} \mathbf {X} ^ {s - 1} + \dots + a _ {0},
\]

where \(a_{i} \in \mathfrak{m}\) for \(0 \leqslant i \leqslant s - 1\) .

Note that the product of two distinguished polynomials is a distinguished polynomial.

PROPOSITION 6 (Preparation Theorem). Let \(f \in \mathbf{B}\) be a series whose reduced series is not zero and \(s\) its reduced order. Then there exists a unique ordered pair \((u, \mathbf{F})\) such that \(u\) is an invertible element of \(\mathbf{B}\) , \(\mathbf{F}\) a distinguished polynomial of degree \(s\) and \(f = u\mathbf{F}\) .

We write \(F = X^{s} + G\) , where \(G = g_{0} + \cdots + g_{s-1}X^{s-1} (g_{i} \in A)\) . The relation f = uF is equivalent to \(F = u^{-1}f\) , that is to \(X^{s} = u^{-1}f - G\) . Hence Proposition 5 shows the uniqueness of G and \(u^{-1}\) and therefore of F and u. It also shows that there exist \(v \in B\) and a polynomial \(G = g_{0} + \cdots + g_{s-1}X^{s-1} (g_{i} \in A)\) such that \(X^{s} = vf - G\) ; it remains to show that v is invertible in B and that \(g_{i} \in m\) for all i. Now, writing \(\bar{g}_{i}\) for the canonical image of \(g_{i}\) in k and \(\bar{f}\) , \(\bar{v}\) for the reduced series of f, g,

\[
\mathbf {X} ^ {s} + \bar {g} _ {0} + \bar {g} _ {1} \mathbf {X} + \dots + \bar {g} _ {s - 1} \mathbf {X} ^ {s - 1} = \bar {f} \bar {v};
\]

since \(\bar{f}\) is of order \(s\) , \(\bar{g}_i = 0\) for all \(i\) and \(\bar{v}\) is of order 0, hence \(v\) is invertible.

PROPOSITION 7. Let F be a distinguished polynomial and g, h two formal power series of B such that F = gh. Then there exists an invertible element u of B such that ug and \(u^{-1}h\) are distinguished polynomials and \(\mathbf{F} = (ug)(u^{-1}h)\) .

In fact, the reduced series of g and h are \(\neq0\) ; hence, by Proposition 6, there exist invertible elements u, v of B such that ug and vh are distinguished polynomials. Then \(uvF = (ug)(vh)\) is a distinguished polynomial and uv is invertible. Passing to the reduced series, it is seen immediately that F and uvF have the same reduced order, that is the same degree. The uniqueness assertion in Proposition 6 therefore shows that F = uvF, whence uv = 1.

COROLLARY. Suppose further that A is an integral domain and F a distinguished polynomial of degree s. For F to be extremal in A{[}X{]}, it is necessary and sufficient that it be extremal in B = A{[}{[}X{]}{]}.

Suppose that F is not extremal in A{[}X{]}, so that \(F = f_{1}f_{2}\) , where \(f_{1}\) and \(f_{2}\) are non-invertible elements of A{[}X{]}; the product of the dominant coefficients of \(f_{1}\) and \(f_{2}\) being equal to 1, these coefficients are invertible in A and the hypothesis implies that \(f_{1}\) and \(f_{2}\) are of degrees \textgreater0 and \textless s; as the reduced polynomials \(\bar{f}_{1}, \bar{f}_{2}\) satisfy \(\bar{f}_{1}\bar{f}_{2} = X^{s}\) , neither \(\bar{f}_{1}\) nor \(\bar{f}_{2}\) can be invertible in k{[}{[}X{]}{]}, for, if \(\bar{f}_{1}\) were invertible, \(\bar{f}_{2}\) would be of order s, which is absurd. A fortiori, neither \(f_{1}\) nor \(f_{2}\) is invertible in B and F is not extremal in B.

Conversely, if F is not extremal in A{[}{[}X{]}{]}, then F = gh, where neither g nor h is invertible in B; their reduced orders are therefore \(\geqslant 1\) ; then the distinguished polynomials ug and \(u^{-1}h\) of Proposition 7 are not constant, which shows that F is not extremal in A{[}X{]}.

\subsection{Factoriality of rings of formal power series}

PROPOSITION 8. Let C be a ring which is either a field or a discrete valuation ring. Then the domain of formal power series C{[}{[}X₁, . . ., Xₙ{]}{]} is factorial.

Let p be the maximal ideal of C and \(\pi\) a generator of p (if C is a field, then \(\pi = 0\) ). Let C be given the p-adic topology, which is Hausdorff. As C is a Noetherian local ring, \(B = C[[X_{1}, \ldots, X_{n}]]\) is a Noetherian local ring and its completion is \(\hat{C}[[X_{1}, \ldots, X_{n}]]\) (Chapter III, §2, no. 6, Proposition 6). By the Corollary to Proposition 4 (no. 6), it suffices to prove that \(\hat{C}[[X_{1}, \ldots, X_{n}]]\) is factorial. Now, if C is a field, then \(\hat{C} = C\) ; if C is a discrete valuation ring, the same is true of \(\hat{C}\) (Chapter VI, §5, no. 3, Proposition 5). We shall therefore assume in the remainder of the proof that C is complete.

Arguing by induction starting with the trivial case n = 0, we shall assume that it has been proved that \(A = C[[X_{1}, \ldots, X_{n-1}]]\) is factorial. We shall identify B with \(A[[X_{n}]]\) and denote by m the maximal ideal of A (generated by \(\pi\) , \(X_{1}, \ldots, X_{n-1}\) ). We shall prove that every non-zero element g of B is, in an essentially unique way, a product of extremal elements.

Let K be the field C/Cπ; as B/Bπ is identified with K{[}{[}X₁, . . . , Xₙ{]}{]}, the ideal Bπ is prime and π is extremal. If π ≠ 0, \(B_{B\pi}\) is therefore the ring of a normed discrete valuation w (Chapter VI, § 3, no. 6, Proposition 9); every nonzero element g of B may therefore be written as g = πw(g)f, where f ∈ B and f is not a multiple of π. It will therefore suffice to show that f is an essentially unique product of extremal elements. Now the canonical image of f in K{[}{[}X₁, . . . , Xₙ{]}{]} is not zero; Lemma 3 (no. 7) therefore shows that there exists an automorphism of B which maps f to an element f' such that the coefficients of f'(0, . . . , 0, Xₙ) are not all in Cπ; this means that the coefficients of the series f', considered as a formal power series in Xₙ, are not all in m. It will suffice to prove our assertion for f'.

In what follows, all the elements of B will be considered as formal power series in \(X_{n}\) with coefficients in A. By Proposition 6 of no. 8 (applicable since C and therefore A are separable and complete and the reduced series of \(f'\) is \(\neq0\) ), \(f'\) is associated, in B, with a unique distinguished polynomial F. By Proposition 7 of no. 8, every series which divides \(f'\) (or, what amounts to the same, which divides F) is associated with a distinguished polynomial which divides F and every decomposition of \(f'\) is, to within invertible factors, of the form \(f' = uF_{1}\ldots F_{q}\) , where u is invertible and the \(F_{t}\) are extremal distinguished polynomials (in B) such that \(F = F_{1}\ldots F_{q}\) . By the Corollary to Proposition 7 of no. 8, the \(F_{t}\) are also extremal in \(A[X_{n}]\) . Now, as A is factorial by the induction hypothesis, so is \(A[X_{n}]\) (Theorem 2, no. 5); hence, since they are monic, the \(F_{i}\) are uniquely determined by F (up to a permutation). This shows the uniqueness of the decomposition \(f' = uF_{1} \ldots F_{q}\) ; its existence follows from the fact that B is Noetherian, which completes the proof.

\textbf{Remarks}\par

\begin{enumerate}
\def\labelenumi{(\arabic{enumi})}
\item
  There exist factorial rings A such that the ring A{[}{[}X{]}{]} is not factorial (Exercise 8). However, if A is a principal ideal domain, A{[}{[}X₁, \ldots, Xₙ{]}{]} is factorial (Exercise 9).
\item
  * We shall see later, by homological methods, that every regular local ring is factorial (cf.~§ 4, no. 7, Corollary 3 to Proposition 16). This will give another proof, conceptually simpler, of Proposition 8. *
\end{enumerate}

\section{MODULES OVER INTEGRALLY CLOSED NOETHERIAN DOMAINS}

Throughout this paragraph, A will be a commutative integral domain with field of fractions K. Starting with no. 2, A will be further assumed to be Noetherian and integrally closed (and hence a Krull domain (§ 1, no. 3, Corollary to Theorem 2)); then P(A), D(A) and C(A) will respectively denote the set of prime ideals of A of height 1 (§ 1, no. 6), the divisor group of A (§ 1, no. 3) and the divisor class group of A (§ 1, no. 10), these latter being written additively.

The general method of studying finitely generated modules over an integrally closed Noetherian domain A consists of ``localizing'' the modules with respect to all the prime ideals \(p \in \mathrm{P}(\mathrm{A})\) of height 1 in A; as \(A_{p}\) is then a discrete valuation ring (§ 1, no. 6, Theorem 4), the structure of finitely generated \(A_{p}\) -modules is well known (Algebra, Chapter VII, § 4) and therefore gives information about the structure of finitely generated A-modules. In the particular case where A is a Dedekind domain, we can arrive at as complete a theory as when A is a principal ideal domain (no. 10).

\subsection{Lattices}

DEFINITION 1. Let V be a finite-dimensional vector space over the field K. A lattice of V with respect to A (or simply a lattice of V) is defined to be any sub-A-module M of V satisfying the following condition:

There exist two free sub-A-modules \(L_{1}\) , \(L_{2}\) of V such that \(L_{1} \subset M \subset L_{2}\) and \(\operatorname{rg}_{\mathbf{A}}(\mathbf{L}_{1}) = \operatorname{rg}_{\mathbf{K}}(\mathbf{V})\) .

\textbf{Examples}\par

\begin{enumerate}
\def\labelenumi{(\arabic{enumi})}
\item
  If we take \(\mathbf{V} = \mathbf{K}\) , the lattices of \(\mathbf{K}\) are just the fractional ideals \(\neq (0)\) of \(\mathbf{K}\) ( \(\S 1\) , no. 1, Definition 1).
\item
  If \(\operatorname{rg}_{\mathbf{K}}(\mathbf{V}) = n\) , every free sub-A-module L of V has a basis containing at most n elements, every subset of V which is free over A being free over K; for L to be a lattice of V, it is necessary and sufficient that L have a basis of n elements (in other words, that \(\operatorname{rg}_{\mathbf{A}}(\mathbf{L}) = n\) ).
\item
  If A is a principal ideal domain, every lattice M of V is a finitely generated A-module (since A is Noetherian) which is torsion-free and hence a free A-module (Algebra, Chapter VII, § 4, no. 3, Corollary 2 to Theorem 2).
\end{enumerate}

PROPOSITION 1. For a sub-A-module M of V to be a lattice of V, it is necessary and sufficient that KM = V and that M be contained in a finitely generated sub-A-module of V.

The conditions are obviously necessary, for a free sub-A-module of V with the same rank as V generates V. Conversely, if KM = V, M contains a basis \((a_{i})_{1 \leqslant i \leqslant n}\) of V over K and hence it contains the free sub-A-module \(L_{1}\) generated by the \(a_{i}\) ; on the other hand, if \(M \subset M_{1}\) , where \(M_{1}\) is a sub-A-module of V generated by a finite number of elements \(b_{j}\) and \((e_{i})_{1 \leqslant i \leqslant n}\) is a basis of V over K, there exists an element \(s \neq 0\) of A such that each of the \(b_{j}\) is a linear combination of the \(s^{-1}e_{i}\) with coefficients in A; if \(L_{2}\) is the free sub-A-modules of V generated by the \(s^{-1}e_{i}\) , then \(M \subset L_{2}\) .

COROLLARY. Suppose that A is Noetherian; for a sub-A-module M of V to be a lattice of V, it is necessary and sufficient that KM = V and M be finitely generated.

Remark (1). Recall that for every sub-A-module M of V the canonical mapping \(M \otimes_{A} K \to V\) is injective and has image KM (Algebra, Chapter II, § 7, no. 10, Proposition 26); to say that KM = V means therefore that this mapping is bijective.

PROPOSITION 2. Let M be a lattice of V and \(M_{1}\) a sub-A-module of V. If there exist two elements x, y of \(K^{*}\) such that \(xM \subset M_{1} \subset yM\) , \(M_{1}\) is a lattice of V; conversely, if \(M_{1}\) is a lattice of V, there exist two non-zero elements a, b of A such that \(aM \subset M_{1} \subset b^{-1}M\) .

If \(L_{1}\) , \(L_{2}\) are two free lattices of V such that \(L_{1} \subset M \subset L_{2}\) , the relations \(xM \subset M_{1} \subset yM\) imply \(xL_{1} \subset M_{1} \subset yL_{2}\) and \(xL_{1}\) and \(yL_{2}\) are free lattices; conversely, if \(M_{1}\) is a lattice and \((e_{i})_{1 \leqslant i \leqslant n}\) a basis of \(L_{2}\) over A, the relation \(KM_{1} = V\) implies the existence of \(x = a/s \in K^{*}\) (where a and s are non-zero elements of A) such that \(xe_{i} \in M_{1}\) for all i, whence \(xM \subset xL_{2} \subset M_{1}\) and a fortiori \(aM \subset M_{1}\) ; exchanging the roles of M and \(M_{1}\) it can be similarly shown that there exists \(b \neq 0\) in A such that \(bM_{1} \subset M\) .

PROPOSITION 3. (i) If \(\mathbf{M}_1\) and \(\mathbf{M}_2\) are lattices of \(\mathbf{V}\) , so are \(\mathbf{M}_1 \cap \mathbf{M}_2\) and \(\mathbf{M}_1 + \mathbf{M}_2\) .

\begin{enumerate}
\def\labelenumi{(\roman{enumi})}
\setcounter{enumi}{1}
\item
  If \(\mathbf{W}\) is a vector subspace of \(\mathbf{V}\) and \(\mathbf{M}\) is a lattice of \(\mathbf{V}\) , \(\mathbf{M} \cap \mathbf{W}\) is a lattice of \(\mathbf{W}\) .
\item
  Let \(\mathbf{V},\mathbf{V}_1,\ldots ,\mathbf{V}_k\) be vector spaces of finite rank over \(\mathbf{K}\) and let
\end{enumerate}

\[
f \colon \mathrm{V} _ {1} \times \dots \times \mathrm{V} _ {k} \rightarrow \mathrm{V}
\]

be a multilinear mapping whose image generates V. If \(M_{i}\) is a lattice of \(V_{i}\) for \(1 \leqslant i \leqslant k\) , the sub-A-module of V generated by \(f(\mathbf{M}_{1} \times \cdots \times \mathbf{M}_{k})\) is a lattice of V.

\begin{enumerate}
\def\labelenumi{(\roman{enumi})}
\setcounter{enumi}{3}
\item
  Let \(\mathbf{V}\) and \(\mathbf{W}\) be two vector spaces of finite rank over \(\mathbf{K}\) , \(\mathbf{M}\) a lattice of \(\mathbf{V}\) and \(\mathbf{N}\) a lattice of \(\mathbf{W}\) . The sub-A-module \(\mathbf{N}:\mathbf{M}\) of \(\operatorname{Hom}_{\mathbf{K}}(\mathbf{V},\mathbf{W})\) , consisting of the K-linear mappings \(f\) such that \(f(\mathbf{M})\subset \mathbf{N}\) , is a lattice of \(\operatorname{Hom}_{\mathbf{K}}(\mathbf{V},\mathbf{W})\) .
\item
  By virtue of Proposition 2, there exist non-zero \(a\) and \(b\) in \(\mathbf{A}\) such that \(a\mathbf{M}_1 \subset \mathbf{M}_2 \subset b^{-1}\mathbf{M}_1\) ; we conclude that \(\mathbf{M}_1 \cap \mathbf{M}_2\) and \(\mathbf{M}_1 + \mathbf{M}_2\) lie between \(a\mathbf{M}_1\) and \(b^{-1}\mathbf{M}_1\) and are therefore lattices by virtue of Proposition 2.
\item
  Let S be a complement of W in V, \(L_{w}\) a free lattice of W and \(L_{s}\) a free lattice of S, so that \(L = L_{w} \oplus L_{s}\) is a free lattice of V. Then there exist x, y in \(K^{*}\) such that \(xL \subset M \subset yL\) . We deduce that \(xL_{w} \subset M \cap W \subset yL_{w}\) , which shows that \(M \cap W\) is a lattice of W (Proposition 2).
\item
  As \(KM_{i} = V_{i}\) , clearly by linearity \(f(\mathbf{M}_{1} \times \cdots \times \mathbf{M}_{k})\) generates the vector K-space V; on the other hand, for all i, there exists a finitely generated sub-A-module \(N_{i}\) of \(V_{i}\) such that \(M_{i} \subset N_{i}\) ; the sub-A-module N of V generated by \(f(\mathbf{N}_{1} \times \cdots \times \mathbf{N}_{k})\) is finitely generated and contains M and hence M is a lattice of V (Proposition 1).
\item
  Let P (resp. Q) be a free lattice of V (resp. W) containing M (resp. contained in N); obviously N: M ⊃ Q: P. Now it is immediate that Q: P is isomorphic to Hom \(_{A}\) (P, Q), hence is a free A-module of rank (rg \(_{A}\) P)(rg \(_{A}\) Q) (Algebra, Chapter II, § 1, no. 6, Corollary 1 to Proposition 6) and therefore a lattice of Hom \(_{K}\) (V, W). Similarly, if P' (resp. Q') is a free lattice of V (resp. W) contained in M (resp. containing N), then Q': P' ⊃ N: M and Q': P' is a lattice of Hom \(_{K}\) (V, W); whence the conclusion.
\end{enumerate}

\textbf{Remarks}\par

\begin{enumerate}
\def\labelenumi{(\arabic{enumi})}
\setcounter{enumi}{1}
\item
  Proposition 3 (i) shows that the set R(V) of lattices of V is lattice-ordered with respect to inclusion; moreover, if M is a fixed lattice of V, the xM, where x runs through K*, form a subset of R(V) which is both coinitial and cofinal (Set Theory, Chapter III, § 1, no. 7).
\item
  In the notation of Proposition 3 (iv), the canonical mapping
\end{enumerate}

\[
\mathrm{N}: \mathrm{M} \rightarrow \operatorname{Hom} _ {\mathrm{A}} (\mathrm{M}, \mathrm{N}),
\]

which maps every K-linear mapping \(f \in N\) : M to the A-linear mapping from M to N which has the same graph as \(f|_{M}\) , is bijective; for every A-linear mapping \(g\) : \(M \to N\) can be imbedded in a K-linear mapping

\[
g \otimes 1: M \otimes_ {A} K \rightarrow N \otimes_ {A} K
\]

and it has been seen that \(M \otimes_{A} K\) and \(N \otimes_{A} K\) are respectively identified with V and W.

In particular, if we take W = K, N = A, \(\operatorname{Hom}_{K}(V, W)\) is just the dual vector K-space \(V^{*}\) of V and A: M is identified with the dual A-module \(M^{*}\) of M; we shall henceforth make this identification and we shall say that \(M^{*}\) is the dual lattice of M: it is therefore the set of \(x^{*} \in V^{*}\) such that \(\langle x, x^{*} \rangle \in A\) for all \(x \in M\) .

COROLLARY. Let U, V, W be three vector spaces of finite rank over K and \(f: \mathbf{U} \times \mathbf{V} \to \mathbf{W}\) a left non-degenerate K-bilinear mapping (Algebra, Chapter IX, § 1, no. 1, Definition 3). If M is a lattice of V and N a lattice of W, the set N: \(f\mathbf{M}\) of \(x \in \mathbf{U}\) such that \(f(x, y) \in \mathbf{N}\) for all \(y \in \mathbf{M}\) is a lattice of U.

Let \(s_f \colon \mathbf{U} \to \operatorname{Hom}_{\mathbf{K}}(\mathbf{V}, \mathbf{W})\) be the K-linear mapping left associated with \(f\) (Algebra, Chapter IX, loc. cit.) such that \(s_f(x)\) is the linear mapping \(y \mapsto f(x, y)\) ; recall that to say that \(f\) is left non-degenerate means that \(s_f\) is injective. By Proposition 3 (iv), \(\mathbf{N} \colon \mathbf{M}\) is a lattice of \(\operatorname{Hom}_{\mathbf{K}}(\mathbf{V}, \mathbf{W})\) ; as \(\mathbf{N} \colon {}_f \mathbf{M} = s_f^{-1}(\mathbf{N} \colon \mathbf{M})\) and \(s_f\) is injective, the corollary follows from Proposition 3 (ii).

\textbf{Examples}\par

\begin{enumerate}
\def\labelenumi{(\arabic{enumi})}
\setcounter{enumi}{3}
\tightlist
\item
  Let S be a (not necessarily associative) K-algebra of finite rank with a unit element; then the bilinear mapping \((x,y)\mapsto xy\) of S × S to S is (left and right) non-degenerate. If M and N are lattices of S with respect to A, so are M.N (Proposition 3 (iii)) and the set of \(x\in S\) such that \(xM\subset N\) (Corollary to Proposition 3). Note that there exists a sub-A-algebra of S containing the unit element of S which is a lattice of S; for consider a basis \((e_{i})_{1\leqslant i\leqslant n}\) of S such that \(e_{1}\) is the unit element of S and let \(e_{i}e_{j}=\sum_{k}c_{ijk}e_{k}\) be the multiplication table of S ( \(1\leqslant i\leqslant n,1\leqslant j\leqslant n\) ), so that \(c_{1jk}=\delta_{jk},c_{11k}=\delta_{1k}\) (Kronecker symbols). Let \(s\in A\) be a non-zero and such that \(c_{ijk}^{\prime}=s.c_{ijk}\in A\) for all triplets of indices \((i,j,k)\) ; if we write \(e_{i}^{\prime}=s^{-1}e_{i}\) for \(i\geqslant2\) , then
\end{enumerate}

\[
e _ {i} ^ {\prime} e _ {j} ^ {\prime} = s c _ {i j 1} ^ {\prime} e _ {1} + \sum_ {k \geqslant 2} c _ {i j k} ^ {\prime} e _ {k} ^ {\prime}
\]

for \(i \geqslant 2\) and \(j \geqslant 2\) ; the lattice of S with basis \(e_1\) and the \(e_i' (2 \leqslant i \leqslant n)\) is a sub-A-algebra of S with unit element \(e_1\) .

\begin{enumerate}
\def\labelenumi{(\arabic{enumi})}
\setcounter{enumi}{4}
\tightlist
\item
  Let V be a finite-dimensional vector space over K and f a non-degenerate bilinear form on V. If M is a lattice of V, it follows from the Corollary to Proposition 3 that the set \(M_{f}^{*}\) of \(x \in V\) such that \(f(x, y) \in A\) for all \(y \in M\) is also a lattice of V; if \(s_{f}: V \to V^{*}\) is the linear mapping left associated with f (which is bijective), \(s_{f}(M_{f}^{*})\) is just the dual lattice \(M^{*}\) of M.
\end{enumerate}

PROPOSITION 4. Let B be an integral domain, A a subring of B and K and L the respective fields of fractions of A and B. Let V be a finite-dimensional vector space over K.

\begin{enumerate}
\def\labelenumi{(\roman{enumi})}
\item
  For every lattice \(\mathbf{M}\) of \(\mathbf{V}\) with respect to \(\mathbf{A}\) , the image \(\mathbf{BM}\) of \(\mathbf{M}_{(\mathbf{B})} = \mathbf{M} \otimes_{\mathbf{A}} \mathbf{B}\) in \(\mathbf{V}_{(\mathbf{L})} = \mathbf{V} \otimes_{\mathbf{K}} \mathbf{L}\) is a lattice of \(\mathbf{V}_{(\mathbf{L})}\) with respect to \(\mathbf{B}\) .
\item
  Suppose further that B is a flat A-module. Then the canonical mapping \(\mathbf{M}_{(\mathbf{B})}\to\mathbf{BM}\) is bijective. If further B is faithfully flat, the mapping which maps every lattice M of V with respect to A to the lattice BM of \(\mathbf{V}_{(\mathbf{L})}\) with respect to B is injective.
\item
  As \(\mathbf{KM} = \mathbf{V}\) , clearly \(\mathbf{L}\) . (BM) = \(\mathbf{V}_{(L)}\) ; on the other hand \(\mathbf{M}\) is contained in a finitely generated sub-A-module \(\mathbf{M}_1\) of \(\mathbf{V}\) and hence \(\mathbf{BM}\) is contained in \(\mathbf{BM}_1\) which is a finitely generated B-module; whence assertion (i) (Proposition 1).
\item
  \(V_{(L)} = V \otimes_{K} L = V \otimes_{A} L\) (Chapter II, § 2, no. 7, Proposition 18) and, as L is a flat B-module, it is also a flat A-module (Chapter I, § 2, no. 7, Corollary 3 to Proposition 8). As B is a flat A-module, the canonical mapping \(M \otimes_{A} B \to V \otimes_{A} B\) is injective; on the other hand, since V is a free K-module and K a flat A-module, V is a flat A-module (Chapter I, § 2, no. 7, Corollary 3 to Proposition 8) and hence the canonical mapping \(V \otimes_{A} B \to V \otimes_{A} L\) is injective, which establishes the first assertion. To see also that the relation \(BM_{1} = BM_{2}\) implies \(M_{1} = M_{2}\) for two lattices \(M_{1}, M_{2}\) of V with respect to A when B is a faithfully flat A-module, note first that \(BM_{1} \cap BM_{2} = B(M_{1} \cap M_{2})\) (Chapter I, § 2, no. 6, Proposition 6); we may therefore confine our attention to the case where \(M_{1} \subset M_{2}\) and our assertion then follows from Chapter I, § 3, no. 1, Proposition 3 applied to the canonical injection \(M_{1} \to M_{2}\) .
\end{enumerate}

COROLLARY. Suppose that A is a discrete valuation ring. Let \(\hat{\mathbf{A}}\) be its completion and let \(\hat{\mathbf{K}}\) be the field of fractions of \(\hat{\mathbf{A}}\) (Chapter VI, § 5, no. 3). The mapping \(\phi\) , which maps every lattice M of V to the lattice \(\hat{\mathbf{A}}\mathbf{M}\) of \(\hat{\mathbf{V}} = \mathbf{V} \otimes_{\mathbf{K}} \hat{\mathbf{K}}\) with respect to \(\hat{\mathbf{A}}\) , is bijective and its inverse maps every lattice M' of \(\hat{\mathbf{V}}\) with respect to \(\hat{\mathbf{A}}\) to its intersection M' \(\cap\) V (V being canonically identified with a vector sub-K-space of \(\hat{\mathbf{V}}\) ).

If L is a free lattice of V, the lattices aL (for \(a \in A\) , \(a \neq 0\) ) form a fundamental system of neighbourhoods of 0 for a topology T on V (compatible with its A-module structure), which (when a basis of L over A is taken) is identified with the product topology on \(K^{n}\) ; by virtue of Proposition 2, a fundamental system of neighbourhoods of 0 for T also consists of all the lattices of V with respect to A; clearly \(\hat{V}\) is the completion of V with respect to T. Moreover, if m is the maximal ideal of A, the topology T induces on every lattice M of V with respect to A the m-adic topology since M is a finitely generated A-module (Chapter III, §3, no. 2, Theorem 2) and \(\hat{A}M\) is the completion of M with respect to this topology (Chapter III, §2, no. 12, Proposition 16); moreover, as M is open (and therefore closed) in V, \(\hat{A}M \cap V = M\) , which proves again the fact that \(\phi\) is injective (which follows directly from Proposition 4, (ii), since \(\hat{A}\) is a faithfully flat A-module). Finally, if \(M'\) is a lattice of \(\hat{V}\) with respect to \(\hat{A}\) , \(M = M' \cap V\) is a lattice of V with respect to A, for every element of \(\hat{A}\) is the product of an element of A and an invertible element of \(\hat{A}\) and hence it follows from Proposition

2 that there exist \(a, b\) in \(\mathbf{A} - \{0\}\) such that \(a\hat{\mathbf{A}}\mathbf{L} \subset \mathbf{M}' \subset b\hat{\mathbf{A}}\mathbf{L}\) , whence \(a\mathbf{L} \subset \mathbf{M}' \cap \mathbf{V} \subset b\mathbf{L}\) . Moreover \(\mathbf{M}'\) is open in \(\mathbf{V}\) and, as \(\mathbf{V}\) is dense in \(\hat{\mathbf{V}}\) , \(\mathbf{M}'\) is the completion of \(\mathbf{M}' \cap \mathbf{V} = \mathbf{M}\) ; this proves that \(\phi\) is surjective, whence the corollary.

Example (6) Let S be a multiplicative subset of A not containing 0; we apply Proposition 4 to \(B = S^{-1}A\) ; then L = K, \(BM = S^{-1}M\) ; hence \(S^{-1}M\) is a lattice of V with respect to \(S^{-1}A\) . Moreover:

PROPOSITION 5. Let V, W vector spaces of finite rank over K, M a lattice of V and N a lattice of W. If M is finitely generated, then (in the notation of Proposition 3):

\[
\mathbf {S} ^ {- 1} (\mathbf {N}: \mathbf {M}) = \mathbf {S} ^ {- 1} \mathbf {N}: \mathbf {S} ^ {- 1} \mathbf {M}\tag{1}
\]

in \(\operatorname{Hom}_{\mathbf{K}}(\mathbf{V},\mathbf{W})\)

Clearly the left hand side of (1) is contained in the right hand side. Conversely, let \(f \in S^{-1}N: S^{-1}M\) and let \((x_i)_{1 \leqslant i \leqslant n}\) be a system of generators of M. There exists \(s \in S\) such that \(f(x_i) \in s^{-1}N\) for all \(i\) and hence \(sf \in N: M\) , which proves the proposition.

\subsection{Duality; reflexive modules}

Recall that from now on the domain A is assumed to be Noetherian and integrally closed and that P(A) (or simply P) denotes the set of prime ideals of A of height 1. Every lattice with respect to A is a finitely generated A-module (no. 1, Corollary to Proposition 1).

Let V be a vector space of finite rank over K, V* its dual and V** its bidual; we shall identify V and V** by means of the canonical mapping \(c_{V}\) (Algebra, Chapter II, § 7, no. 5, Theorem 6). Let M be a lattice of V; recall that the dual A-module M* of M is canonically identified with the dual lattice of M, the set of \(x^{*} \in V^{*}\) such that \(\langle x, x^{*} \rangle \in A\) for all \(x \in M\) ; the bidual A-module M** of M is therefore a lattice of V which contains M. Moreover M*** = M*, for the relation \(M \subset M**\) implies \((\mathbf{M}^{**})^{*} \subset \mathbf{M}^{*}\) and on the other hand \(M^{*} \subset (\mathbf{M}^{*})^{**}\) by the above (cf.~Set Theory, Chapter III, § 1, no. 5, Proposition 2).

If p is a prime ideal, Proposition 5 applied with N = A gives the relation \((\mathbf{M}^{*})_{\mathfrak{p}} = (\mathbf{M}_{\mathfrak{p}})^{*}\) , which justifies the notation \(M_{p}^{*}\) for both terms.

THEOREM 1. If \(\mathbf{M}\) is a lattice of \(\mathbf{V}\) , then \(\mathbf{M}^{*} = \bigcap_{\mathfrak{p} \in \mathbb{P}} \mathbf{M}_{\mathfrak{p}}^{*}\) .

Clearly \(\mathbf{M}^*\) is contained in each of the \(\mathbf{M}_{\mathfrak{p}}^{*}\) . Conversely, suppose that \(x^{*} \in \bigcap_{\mathfrak{p} \in \mathbf{P}} \mathbf{M}_{\mathfrak{p}}^{*}\) ; if \(x \in \mathbf{M}\) , then \(\langle x, x^{*} \rangle \in \bigcap_{\mathfrak{p} \in \mathbf{P}} \mathbf{A}_{\mathfrak{p}}\) and, as \(\mathbf{A} = \bigcap_{\mathfrak{p} \in \mathbf{P}} \mathbf{A}_{\mathfrak{p}}\) (§ 1, no. 6, Theorem 4), \(x^{*} \in \mathbf{M}^{*}\) .

COROLLARY. \(\mathbf{M}^{**} = \bigcap_{p \in P}\mathbf{M}_{p}\)

Theorem 1 applied to \(\mathbf{M}^*\) shows that \(\mathbf{M}^{**} = \bigcap_{\mathfrak{p}\in \mathbb{P}}\mathbf{M}_{\mathfrak{p}}^{**}\) . But as \(\mathbf{A}_{\mathfrak{p}}\) is a principal ideal domain (§ 1, no. 6, Theorem 4), \(\mathbf{M}_{\mathfrak{p}}\) is a finitely generated free \(\mathbf{A}_{\mathfrak{p}}\) -module and hence \(\mathbf{M}_{\mathfrak{p}}^{**}\) is canonically identified with \(\mathbf{M}_{\mathfrak{p}}\) (Algebra, Chapter II, § 2, no. 7, Proposition 14), whence the corollary.

For any lattice M with respect to A, the canonical mapping \(c_{M}: M \to M^{**}\) (Algebra, Chapter II, §2, no. 7) identifies an element \(x \in M\) with itself, for x is the unique element y of \(V = V^{**}\) such that \(\langle x, x^{*} \rangle = \langle y, x^{*} \rangle\) for all \(x^{*} \in M^{*}\) , since \(M^{*}\) generates \(V^{*}\) . We shall say that M is reflexive if \(M^{**} = M\) (loc. cit.). As we have above \(M^{*} = (M^{*})^{**}\) , it is seen that the dual of any lattice M is always reflexive.

Remark (1) Let M be a finitely generated A-module; it is immediate that the dual M* of M, identified with a sub-A-module of \(\operatorname{Hom}_{\mathbf{A}}(\mathbf{M}, \mathbf{K})\) is a lattice of the vector K-space \(\operatorname{Hom}_{\mathbf{A}}(\mathbf{M}, \mathbf{K})\) ; in particular, every finitely generated reflexive A-module is isomorphic to a lattice of a suitable vector K-space.

THEOREM 2. If M is a lattice of V, the following conditions are equivalent:

\begin{enumerate}
\def\labelenumi{(\alph{enumi})}
\item
  M is reflexive.
\item
  \(\mathbf{M} = \bigcap_{\mathfrak{p}\in \mathbb{P}}\mathbf{M}_{\mathfrak{p}}.\)
\item
  \(\operatorname{Ass}(\mathbf{V} / \mathbf{M})\subset \mathbf{P}\)
\end{enumerate}

The equivalence of (a) and (b) follows from the Corollary to Theorem 1. If (b) holds, V/M is canonically identified with a sub-A-module of the product \(\prod_{\mathfrak{p} \in \mathbb{P}} (\mathrm{V} / \mathrm{M}_{\mathfrak{p}})\) ; but in fact, it is contained in the direct sum \(\bigoplus_{\mathfrak{p} \in \mathbb{P}} (\mathrm{V} / \mathrm{M}_{\mathfrak{p}})\) : for if \(\mathbf{L} \subset \mathbf{M}\) is a free lattice and \((e_i)_{1 \leq i \leq n}\) a basis of \(\mathbf{L}\) , each of the coordinates \(x_i\) of a point \(x \in \mathbf{V}\) with respect to \((e_i)\) belongs to \(\mathbf{A}_{\mathfrak{p}}\) except for a finite number of values of \(\mathfrak{p} \in \mathbb{P}\) . The relation \(\mathrm{V} / \mathrm{M} \subset \bigoplus_{\mathfrak{p} \in \mathbb{P}} (\mathrm{V} / \mathrm{M}_{\mathfrak{p}})\) then implies:

\[
\operatorname{Ass} (\mathrm{V} / \mathrm{M}) \subset \bigcap_ {\mathfrak {p} \in \mathrm{P}} \operatorname{Ass} (\mathrm{V} / \mathrm{M} _ {\mathfrak {p}}).
\]

As \(V / M_{p}\) is an \(A_{p}\) -module, an element of \(A - p\) cannot annihilate an element \(\neq 0\) of \(V / M_{p}\) , since the elements of \(A - p\) are invertible in \(A_{p}\) ; the elements of \(\operatorname{Ass}(V / M_{p})\) are therefore contained in \(p\) and are \(\neq 0\) , since \(V / M_{p}\) is a torsion \(A_{p}\) -module; as \(p\) is of height 1, necessarily \(\operatorname{Ass}(V / M_{p}) = \{p\}\) if \(V / M_{p} \neq \{0\}\) and \(\operatorname{Ass}(V / M_{p}) = \varnothing\) if \(V / M_{p} = \{0\}\) ; hence \(\operatorname{Ass}(V / M) \subset P\) .

Finally, if condition (c) holds, then

\[
\operatorname{Ass} (\mathbf {M} ^ {* *} / \mathbf {M}) \subset \operatorname{Ass} (\mathbf {V} / \mathbf {M}) \subset \mathbf {P}.
\]

On the other hand, if \(\mathfrak{p} \in \mathbf{P}\) , then it has been seen in the proof of the Corollary to Theorem 1 that \(\mathbf{M}_{\mathfrak{p}}^{**} = \mathbf{M}_{\mathfrak{p}}\) , whence \(\mathfrak{p} \notin \operatorname{Ass}(\mathbf{M}^{**}/\mathbf{M})\) (Chapter IV, § 1, no. 3, Corollary 1 to Proposition 7). We conclude that \(\operatorname{Ass}(\mathbf{M}^{**}/\mathbf{M}) = \varnothing\) , whence \(\mathbf{M}^{**} = \mathbf{M}\) (Chapter IV, § 1, no. 1, Corollary 1 to Proposition 2).

COROLLARY. Let M, N be two lattices of V with respect to A such that N is reflexive. In order that \(M \subset N\) , it is necessary and sufficient that, for all \(p \in P\) , \(M_{p} \subset N_{p}\) .

The condition is obviously necessary and, if it is fulfilled, then

\[
\bigcap_ {p \in P} M _ {p} \subset \bigcap_ {p \in P} N _ {p} = N.
\]

As \(\mathbf{M} \subset \mathbf{M}^{**} = \bigcap_{\mathfrak{p} \in \mathbb{P}} \mathbf{M}_{\mathfrak{p}}\) , certainly \(\mathbf{M} \subset \mathbf{N}\) .

Examples

\begin{enumerate}
\def\labelenumi{(\arabic{enumi})}
\item
  Every free lattice is reflexive.
\item
  Take \(\mathbf{V} = \mathbf{K}\) . For a fractional ideal \(\mathfrak{a}\) of \(\mathbf{K}\) to be a reflexive lattice, it is necessary and sufficient that it be a divisorial ideal by virtue of criterion (b) of Theorem 2 and §1, no. 4, Propositions 5 and 7.
\item
  Let M be a lattice with respect to A; if S is a multiplicative subset of A not containing 0, Proposition 5 of no. 1 shows that \(\mathrm{S}^{-1}(\mathbf{M}^{*}) = (\mathrm{S}^{-1}\mathrm{M})^{*}\) ; if M is reflexive, \(S^{-1}M\) is therefore a reflexive lattice with respect to \(S^{-1}A\) .
\end{enumerate}

PROPOSITION 6. (i) If \(\mathbf{M}_1\) and \(\mathbf{M}_2\) are reflexive lattices of \(\mathbf{V}\) , so is \(\mathbf{M}_1 \cap \mathbf{M}_2\) .

\begin{enumerate}
\def\labelenumi{(\roman{enumi})}
\setcounter{enumi}{1}
\item
  If \(\mathbf{W}\) is a vector subspace of \(\mathbf{V}\) and \(\mathbf{M}\) is a reflexive lattice of \(\mathbf{V}\) , \(\mathbf{M} \cap \mathbf{W}\) is a reflexive lattice of \(\mathbf{W}\) .
\item
  Let V, W be two vector spaces of finite rank over K and M (resp. N) a lattice of V (resp. W). If N is reflexive, the lattice N: M of \(\operatorname{Hom}_{\mathbb{K}}(\mathrm{V}, \mathrm{W})\) (no. 1, Proposition 3) is reflexive.
\item
  \((\mathbf{M}_1 \cap \mathbf{M}_2)_p = (\mathbf{M}_1)_p \cap (\mathbf{M}_2)_p\) for all \(p \in P\) (Chapter II, § 2, no. 4, Theorem 1). If \(\mathbf{M}_1 = \bigcap_{p \in P} (\mathbf{M}_1)_p\) and \(\mathbf{M}_2 = \bigcap_{p \in P} (\mathbf{M}_2)_p\) , then
\end{enumerate}

\[
\mathrm{M} _ {1} \cap \mathrm{M} _ {2} = \bigcap_ {\mathfrak {p} \in \mathrm{P}} (\mathrm{M} _ {1} \cap \mathrm{M} _ {2}) _ {\mathfrak {p}}
\]

whence the conclusion by virtue of Theorem 2.

\begin{enumerate}
\def\labelenumi{(\roman{enumi})}
\setcounter{enumi}{1}
\tightlist
\item
  Similarly \((\mathbf{M} \cap \mathbf{W})_{\mathfrak{p}} = \mathbf{M}_{\mathfrak{p}} \cap \mathbf{W}_{\mathfrak{p}} = \mathbf{M}_{\mathfrak{p}} \cap \mathbf{W}\) , whence
\end{enumerate}

\[
\mathbf {M} \cap \mathbf {W} = \bigcap_ {\mathfrak {p} \in \mathbf {P}} (\mathbf {M} \cap \mathbf {W}) _ {\mathfrak {p}},
\]

which proves (ii).

\begin{enumerate}
\def\labelenumi{(\roman{enumi})}
\setcounter{enumi}{2}
\tightlist
\item
  As \(\mathbf{M}\) is finitely generated, it follows from no. 1, Proposition 5 that \((\mathbf{N}:\mathbf{M})_{\mathfrak{p}} = \mathbf{N}_{\mathfrak{p}}:\mathbf{M}_{\mathfrak{p}}\) ; moreover, the relation \(\mathbf{N} = \bigcap_{\mathfrak{p}\in P}\mathbf{N}_{\mathfrak{p}}\) implies:
\end{enumerate}

\[
\mathbf {N}: \mathbf {M} = \bigcap_ {\mathfrak {p} \in \mathbf {P}} \left(\mathbf {N} _ {\mathfrak {p}}: \mathbf {M} _ {\mathfrak {p}}\right).
\]

For if \(f \in \bigcap_{\mathfrak{p} \in \mathbb{P}} (\mathbf{N}_{\mathfrak{p}}: \mathbf{M}_{\mathfrak{p}})\) and \(x \in \mathbf{M}\) , then \(f(x) \in \bigcap_{\mathfrak{p} \in \mathbb{P}} \mathbf{N}_{\mathfrak{p}} = \mathbf{N}\) , whence \(f \in \mathbf{N}: \mathbf{M}\) ; this shows that \(\mathbf{N}: \mathbf{M}\) is reflexive.

\textbf{Remarks}\par

\begin{enumerate}
\def\labelenumi{(\arabic{enumi})}
\setcounter{enumi}{1}
\item
  If \(\mathbf{M}_1\) and \(\mathbf{M}_2\) are reflexive lattices of \(\mathbf{V}\) , the lattice \(\mathbf{M}_1 + \mathbf{M}_2\) is not necessarily reflexive (cf.~§ 1, Exercise 2).
\item
  If M is a finitely generated A-module and T its torsion submodule, the dual \(M^{*}\) of M is the same as the dual of M/T, since, for every linear form f on M, the image \(f(T)\) is a torsion submodule of A and hence zero. As M/T is isomorphic to a lattice of a vector space over K, it is seen that the dual of every finitely generated A-module is reflexive.
\end{enumerate}

PROPOSITION 7. Let \(0 \to M \to N \to Q \to 0\) be an exact sequence of A-modules. Suppose that \(N\) is finitely generated and torsion-free.

\begin{enumerate}
\def\labelenumi{(\roman{enumi})}
\item
  If \(\mathbf{M}\) is reflexive, then \(\operatorname{Ass}(\mathbf{Q}) \subset \mathbf{P} \cup \{\{0\}\}\) (in other words, every ideal associated with \(\mathbf{Q}\) , is, either (0), or of height 1).
\item
  Conversely, if N is reflexive and \(\operatorname{Ass}(\mathbf{Q}) \subset \mathbf{P} \cup \{\{0\}\}\) , then M is reflexive.
\end{enumerate}

As A is Noetherian, M is also finitely generated; if we write \(V = M_{(K)}\) , \(W = N_{(K)}\) , M (resp. N) is canonically identified with a lattice of V (resp. W) (no. 1, Proposition 1). Consider the two exact sequences:

\[
0 \rightarrow \mathrm{V} / \mathrm{M} \rightarrow \mathrm{W} / \mathrm{M} \rightarrow \mathrm{W} / \mathrm{V} \rightarrow 0
\]

\[
0 \rightarrow Q \rightarrow W / M \rightarrow W / N \rightarrow 0.
\]

\begin{enumerate}
\def\labelenumi{(\roman{enumi})}
\tightlist
\item
  We deduce (Chapter IV, § 1, no. 1, Proposition 3) that:
\end{enumerate}

\[
\operatorname{Ass} (Q) \subset \operatorname{Ass} (W / M) \subset \operatorname{Ass} (V / M) \cup \operatorname{Ass} (W / V).
\]

If M is reflexive, then \(\operatorname{Ass}(V/M) \subset P\) (Theorem 2); on the other hand, clearly \(\operatorname{Ass}(W/V)\) is, either empty, or reduced to \(\{0\}\) ; whence (i).

\begin{enumerate}
\def\labelenumi{(\roman{enumi})}
\setcounter{enumi}{1}
\tightlist
\item
  Similarly:
\end{enumerate}

\[
\operatorname{Ass} (\mathbf {V} / \mathbf {M}) \subset \operatorname{Ass} (\mathbf {W} / \mathbf {M}) \subset \operatorname{Ass} (\mathbf {Q}) \cup \operatorname{Ass} (\mathbf {W} / \mathbf {N}).
\]

The hypotheses therefore imply that \(\operatorname{Ass}(\mathbf{V} / \mathbf{M}) \subset \mathbf{P} \cup \{\{0\}\}\) . But \(\mathbf{V} / \mathbf{M}\) is a torsion A-module and hence \(\{0\} \notin \operatorname{Ass}(\mathbf{V} / \mathbf{M})\) ; Theorem 2 then shows that \(\mathbf{M}\) is reflexive.

PROPOSITION 8. Let R and S be two commutative rings, \(\rho: \mathrm{R} \to \mathrm{S}\) a ring homomorphism and M a finitely generated R-module. Suppose that R is Noetherian and that S is a flat R-module. Then, if M is reflexive, so is the S-module \(\mathrm{M}_{(\mathrm{S})} = \mathrm{M} \otimes_{\mathrm{R}} \mathrm{S}\) .

We know (Chapter I, § 2, no. 10, Proposition 11) that there exists a canonical isomorphism \(\omega_{\mathbf{M}}\colon (\mathbf{M}^{*})_{(\mathbf{S})}\to (\mathbf{M}_{(\mathbf{S})})^{*}\) , such that

\[
\langle x \otimes 1, \omega_ {\mathrm{M}} (x ^ {*} \otimes 1) \rangle = \rho (\langle x, x ^ {*} \rangle)
\]

for \(x \in \mathbf{M}\) , \(x^{*} \in \mathbf{M}^{*}\) . As \(\mathbf{M}\) is a quotient of a finitely generated free R-module \(\mathbf{L}\) , \(\mathbf{M}^{*}\) is isomorphic to a sub-R-module of the dual \(\mathbf{L}^{*}\) and \(\mathbf{L}^{*}\) is free and finitely generated; since \(\mathbf{R}\) is Noetherian, \(\mathbf{M}^{*}\) is therefore also a finitely generated R-module, whence an isomorphism \(\omega_{\mathbf{M}^{*}}: (\mathbf{M}^{**})_{(\mathbf{S})} \to ((\mathbf{M}^{*})_{(\mathbf{S})})^{*}\) such that

\[
\langle x ^ {*} \otimes 1, \omega_ {M ^ {*}} (x ^ {* *} \otimes 1) \rangle = \rho (\langle x ^ {*}, x ^ {* *} \rangle)
\]

for \(x^{*} \in \mathbf{M}^{*}\) and \(x^{**} \in \mathbf{M}^{**}\) . On the other hand, there is an isomorphism \(^t\omega_{\mathbf{M}}: (\mathbf{M}_{(\mathbf{S})})^{**} \to ((\mathbf{M}^{*})_{(\mathbf{S})})^{*}\) , whence by composition a canonical isomorphism:

\[
\phi = \left(^ {t} \omega_ {\mathrm{M}} ^ {- 1}\right) \circ (\omega_ {\mathrm{M} ^ {*}}): (\mathrm{M} ^ {* *}) _ {(\mathrm{S})} \rightarrow (\mathrm{M} _ {(\mathrm{S})}) ^ {* *}
\]

such that, in the above notation:

\[
\langle \omega_ {M} (x ^ {*} \otimes 1), \phi (x ^ {* *} \otimes 1) \rangle = \rho (\langle x ^ {*}, x ^ {* *} \rangle).\tag{1}
\]

We consider now the canonical homomorphism \(c_{\mathbf{M}}: \mathbf{M} \to \mathbf{M}^{**}\) and show that the composite homomorphism:

\[
\psi : \mathrm{M} _ {(\mathbf {S})} \xrightarrow [ c _ {\mathbf {M}} \otimes 1 ]{} (\mathrm{M} ^ {* *}) _ {(\mathbf {S})} \longrightarrow (\mathrm{M} _ {(\mathbf {S})}) ^ {* *}\tag{2}
\]

is just the canonical homomorphism \(c_{\mathbf{M}_{(\mathbf{S})}}\) . This follows immediately from (1) which gives the relations:

\[
\begin{array}{r l} \langle \omega_ {\mathrm{M}} (x ^ {*} \otimes 1), \psi (x \otimes 1) \rangle = & \rho (\langle x ^ {*}, c _ {\mathrm{M}} (x) \rangle) \\ & = \rho (\langle x, x ^ {*} \rangle) = \langle x \otimes 1, \omega_ {\mathrm{M}} (x ^ {*} \otimes 1) \rangle \end{array}
\]

and from the fact that the elements \(\omega_{\mathbf{M}}(x^{*}\otimes 1)\) generate \((\mathbf{M}_{(\mathbf{S})})^{*}\) . This being so, the hypothesis that \(\mathbf{M}\) is reflexive means that \(c_{\mathbf{M}}\) is bijective, hence so is \(c_{\mathbf{M}}\otimes 1\) and therefore \(\psi = c_{\mathbf{M}_{(\mathbf{S})}}\) is bijective, which shows the proposition.

\subsection{Local construction of reflexive modules}

We keep the notation and hypotheses of no. 2. We shall say that a property holds ``for almost all \(p \in P\) '' if the set of \(p \in P\) for which it is not true is finite.

THEOREM 3. Let V be a vector space of finite rank over K and M a lattice of V with respect to A.

\begin{enumerate}
\def\labelenumi{(\roman{enumi})}
\item
  Let N be a lattice of V with respect to A; then, for every prime ideal p of A, \(N_{p}\) is a lattice of V with respect to \(A_{p}\) and, for almost all \(p \in P\) , \(N_{p} = M_{p}\) .
\item
  Conversely, suppose given for all \(\mathfrak{p} \in \mathbf{P}\) a lattice \(\mathrm{N}(\mathfrak{p})\) of \(\mathbf{V}\) with respect to \(\mathbf{A}_{\mathfrak{p}}\) such that \(\mathrm{N}(\mathfrak{p}) = \mathrm{M}_{\mathfrak{p}}\) for almost all \(\mathfrak{p} \in \mathbf{P}\) . Then \(\mathrm{N} = \bigcap_{\mathfrak{p} \in \mathbf{P}} \mathrm{N}(\mathfrak{p})\) is a reflexive lattice of \(\mathbf{A}\) with respect to \(\mathbf{A}\) and it is the only reflexive lattice \(\mathrm{N}'\) of \(\mathbf{V}\) with respect to \(\mathbf{A}\) such that \(\mathrm{N}_{\mathfrak{p}}' = \mathrm{N}(\mathfrak{p})\) for all \(\mathfrak{p} \in \mathbf{P}\) .
\item
  The first assertion follows from no. 1, Proposition 4. Moreover, there exist \(x, y\) in \(\mathbf{K}^*\) such that \(x\mathrm{N} \subset \mathrm{M} \subset y\mathrm{N}\) (no. 1, Proposition 2); we know that, for almost all \(\mathfrak{p} \in \mathbf{P}\) , \(v_{\mathfrak{p}}(x) = v_{\mathfrak{p}}(y) = 0\) (§ 1, no. 6, Theorem 4), which shows that \(x\) and \(y\) are invertible in \(\mathbf{A}_{\mathfrak{p}}\) and hence \(\mathbf{M}_{\mathfrak{p}} = \mathbf{N}_{\mathfrak{p}}\) .
\item
  We may replace M by \(x^{-1}\mathbf{M}\) where \(x \neq 0\) in A and assume that \(\mathrm{N}(\mathfrak{p}) \subset \mathbf{M}_{\mathfrak{p}}\) for all \(\mathfrak{p} \in \mathrm{P}\) . Let \(\mathfrak{p}_1, \ldots, \mathfrak{p}_h\) be the elements of P such that \(\mathrm{N}(\mathfrak{p}) = \mathrm{M}_{\mathfrak{p}}\) for \(\mathfrak{p}\) distinct from the \(\mathfrak{p}_i\) ( \(1 \leqslant i \leqslant h\) ); we write:
\end{enumerate}

\[
Q = M \cap N (p _ {1}) \cap \dots \cap N (p _ {h}).
\]

As each of the \(\mathrm{N}(p_{i})\) contain a free lattice with respect to \(A_{p_{i}}\) , it contains a fortiori a lattice of V with respect to A, hence Q contains a lattice of V with respect to A (no. 1, Proposition 3) and, as Q is contained in M, Q is a lattice with respect to A. To prove that \(Q_{p} = N(p)\) for all \(p \in P\) , we shall use the following lemma:

LEMMA 1. Let \(\mathfrak{p}\) and \(\mathfrak{p}'\) be two prime ideals of \(\mathbf{A}\) such that (0) is the only prime ideal of \(\mathbf{A}\) contained in \(\mathfrak{p} \cap \mathfrak{p}'\) . Then, for every sub- \(\mathbf{A}\) -module \(\mathbf{E}\) of \(\mathbf{V}\) , \((\mathbf{E}_{\mathfrak{p}})_{\mathfrak{p}'} = \mathbf{K} \cdot \mathbf{E}\) .

Let S be the multiplicative subset \((\mathbf{A}-\mathfrak{p})(\mathbf{A}-\mathfrak{p}')\) of A; by Chapter II, §2, no. 3, Proposition 7, \((\mathrm{E}_{\mathfrak{p}})_{\mathfrak{p}'}=\mathrm{S}^{-1}\mathrm{E}\) . Further, \(A\subset S^{-1}A\subset K\) ; the prime ideals of \(S^{-1}A\) correspond to the prime ideals q of A such that \(q\cap S=\varnothing\) (Chapter II, §2, no. 5, Proposition 11) and by hypothesis (0) is the only prime ideal of A not meeting S; hence \(S^{-1}A=K\) and \(S^{-1}E=K.E\) .

We now return to the proof of (ii). If \(p \in P\) is distinct from the \(p_i (1 \leqslant i \leqslant h)\) , Lemma 1 applied to \(\mathrm{N}(\mathfrak{p}_i)\) gives \((\mathrm{N}(\mathfrak{p}_i))_{\mathfrak{p}} = ((\mathrm{N}(\mathfrak{p}_i))_{\mathfrak{p}_i})_{\mathfrak{p}} = \mathrm{K}. \mathrm{N}(\mathfrak{p}_i) = \mathrm{V}\) , since the \(p_i\) and p are of height 1. Then

\[
\mathbf {Q} _ {\mathfrak {p}} = \mathbf {M} _ {\mathfrak {p}} \cap (\mathbf {N} (\mathfrak {p} _ {1})) _ {\mathfrak {p}} \cap \dots \cap (\mathbf {N} (\mathfrak {p} _ {h})) _ {\mathfrak {p}} = \mathbf {M} _ {\mathfrak {p}} = \mathbf {N} (\mathfrak {p})
\]

(Chapter II, § 2, no. 4). On the other hand, if \(\mathfrak{p}\) is equal to \(\mathfrak{p}_i\) ( \(1 \leqslant i \leqslant h\) ), then \((\mathrm{N}(\mathfrak{p}_i))_{\mathfrak{p}_j} = \mathrm{V}\) for \(i \neq j\) by the argument as above and \((\mathrm{N}(\mathfrak{p}_i))_{\mathfrak{p}_i} = \mathrm{N}(\mathfrak{p}_i)\) , whence

\[
\mathbf {Q} _ {\mathfrak {p} _ {i}} = \mathbf {M} _ {\mathfrak {p} _ {i}} \cap \mathbf {N} (\mathfrak {p} _ {i}) = \mathbf {N} (\mathfrak {p} _ {i}).
\]

We have therefore proved that \(Q_{\mathfrak{p}} = N(\mathfrak{p})\) for all \(\mathfrak{p} \in P\) . Then

\[
\mathbf {N} = \mathbf {Q} ^ {* *} = \bigcap_ {\mathfrak {p} \in \mathbf {P}} \mathbf {Q} _ {\mathfrak {p}}
\]

is reflexive and satisfies the relations \(N_{p} = Q_{p} = N(p)\) for all \(p \in P\) ; the uniqueness property follows immediately from Theorem 2 of no. 2.

Remark. Let L be a free lattice of V with respect to A. Since \(A_{p}\) is a principal ideal domain for \(p \in P\) , \(N(p)\) is a free \(A_{p}\) -module of the same rank as L and there exists \(u(p) \in \mathbf{GL}(V)\) such that \(u(p)(L_{p}) = N_{p}\) ; this condition moreover determines \(u(p)\) to within right multiplication by an element of \(\mathbf{GL}(L_{p})\) . The condition \(N(p) = L_{p}\) for almost all \(p \in P\) means that necessarily \(u(p) \in \mathbf{GL}(L_{p})\)

for almost all \(\mathfrak{p} \in \mathbb{P}\) . The families \((u(\mathfrak{p}))_{\mathfrak{p} \in \mathbb{P}}\) satisfying the latter property form a multiplicative group \(\mathbf{GL}_a(V)\) containing as subgroup the product \(\prod_{\mathfrak{p} \in \mathbb{P}} \mathbf{GL}(L_{\mathfrak{p}})\) . Theorem 3 then shows that the set of reflexive lattices of \(V\) is in canonical one-to-one correspondence with the homogeneous space \(\mathbf{GL}_a(V)/\prod_{\mathfrak{p} \in \mathbb{P}} \mathbf{GL}(L_{\mathfrak{p}})\) . If a basis \((e_i)_{1 \leq i \leq n}\) of \(L\) over \(A\) is chosen, \(\mathbf{GL}(V)\) (resp. \(\mathbf{GL}(L_{\mathfrak{p}})\) ) is identified with the group of invertible matrices \(\mathbf{GL}(n, K)\) (resp. \(\mathbf{GL}(n, A_{\mathfrak{p}})\) ) and the group \(\mathbf{GL}_a(V)\) with the group of systems of matrices of order \(n\) , \((U(\mathfrak{p}))_{\mathfrak{p} \in P}\) , such that \(U(\mathfrak{p}) \in \mathbf{GL}(n, K)\) for all \(\mathfrak{p} \in P\) and \(U(\mathfrak{p}) \in \mathbf{GL}(n, A_{\mathfrak{p}})\) for almost all \(\mathfrak{p} \in P\) . If \(A\) is a Dedekind domain, the group \(\mathbf{GL}_a(V)\) is also identified with the group \(\mathbf{GL}(n, A)\) , where \(A\) is the ring of restricted adèles (§ 2, no. 4).

\subsection{Pseudo-isomorphisms}

We preserve the notation and hypotheses of nos. 2 and 3.

PROPOSITION 9. Let M be a finitely generated A-module. The following conditions are equivalent:

\begin{enumerate}
\def\labelenumi{(\alph{enumi})}
\item
  \(\mathbf{M}_{\mathfrak{p}} = 0\) for every prime ideal \(\mathfrak{p}\) of height \(\leqslant 1\) .
\item
  The annihilator \(\mathfrak{a}\) of \(\mathbf{M}\) is an ideal \(\neq (0)\) and \(\mathbf{A}:\mathfrak{a} = \mathbf{A}\) ( \(\mathbf{A}:\mathfrak{a}\) denoting, as in § 1, no. 1, the set of \(x\in \mathbf{K}\) such that \(xa\in \mathbf{A}\) ).
\end{enumerate}

We know (Chapter II, § 2, no. 2, Corollary 2 to Proposition 4) that the condition \(\mathbf{M}_{\mathfrak{p}} = 0\) is equivalent to \(\mathfrak{a} \nsubseteq \mathfrak{p}\) and hence to \(\mathfrak{a}\mathbf{A}_{\mathfrak{p}} = \mathbf{A}_{\mathfrak{p}}\) (Chapter II, § 2, no. 5, Remark); on the other hand, for every integral ideal \(\mathfrak{b} \neq 0\) of A, the relation `` \(\mathfrak{b}\mathbf{A}_{\mathfrak{p}} = \mathbf{A}_{\mathfrak{p}}\) for all \(\mathfrak{p} \in \mathbb{P}\) '' is equivalent to \(\operatorname{div} \mathfrak{b} = \operatorname{div} \mathbf{A} = 0\) in D(A) (§ 1, no. 4, Proposition 7), or also to \(\operatorname{div}(\mathbf{A}: \mathfrak{b}) = 0\) and, as A: \(\mathfrak{b}\) is divisorial (§ 1, no. 1, Proposition 1), this relation is also equivalent to A: \(\mathfrak{b} = \mathbf{A}\) . The proposition then follows by noting that to say that \(\mathfrak{a} \nsubseteq \mathfrak{p}\) for \(\mathfrak{p} = (0)\) means that \(\mathfrak{a} \neq (0)\) .

Remark (1) The equivalent conditions of Proposition 9 mean also that \(\operatorname{Ass}(\mathbf{M})\) contains no prime ideal of height \(\leqslant 1\) . * They may also be interpreted by saying that \(\operatorname{Supp}(\mathbf{M})\) is of codimension \(\geqslant 2\) in \(\operatorname{Spec}(\mathbf{A})\) . *

DEFINITION 2. An A-module M is called pseudo-zero if it is finitely generated and it satisfies the equivalent conditions of Proposition 9.

This definition and Proposition 9 show that a pseudo-zero A-module is a torsion A-module; the converse is false.

\textbf{Examples}\par

\begin{enumerate}
\def\labelenumi{(\arabic{enumi})}
\item
  If A is a Dedekind domain, every prime ideal of A is of height \(\leqslant1\) ; to say that M is pseudo-zero means then that Supp(M) = \(\varnothing\) and hence that M = 0 (Chapter II, § 4, no. 4).
\item
  Let \(k\) be a field and \(A = k[X, Y]\) the polynomial ring over \(k\) in two indeterminates; if \(m\) is the maximal ideal \(AX + AY\) of \(A\) , the \(A\) -module \(A/m\) is pseudo-zero; for its annihilator \(m\) is not of height \(\leqslant 1\) since it contains the principal prime ideals \(AX\) and \(AY\) and is distinct from them; therefore \(A: m = A\) (§ 1, no. 6, Corollary 1 to Theorem 3).
\end{enumerate}

DEFINITION 3. Let M and N be two A-modules and f: M → N a homomorphism. f is called pseudo-injective (resp. pseudo-surjective, pseudo-zero) if Ker(f) (resp. Coker(f), Im(f)) is pseudo-zero; f is called pseudo-bijective if it is both pseudo-injective and pseudo-surjective.

A pseudo-bijective homomorphism is also called a pseudo-isomorphism.

Suppose that M and N are finitely generated; then, for \(f: M \to N\) to be pseudo-injective (resp. pseudo-surjective, pseudo-zero), it is necessary and sufficient that, for all \(p \in P \cup \{\{0\}\}, f_{p}: M_{p} \to N_{p}\) be injective (resp. surjective, zero); this follows from the flatness of the A-module \(A_{p}\) (cf.~Chapter I, § 2, no. 3, Remark 2).

Example (3) Let M be a torsion-free finitely generated A-module; then the canonical mapping \(c_{M}: M \to M^{**}\) of M to its bidual is a pseudo-isomorphism. For M is identified with a lattice of \(V = M \otimes_{A} K\) (no. 1, Proposition 1); we have seen that \(M_{p} = M_{p}^{**}\) for all \(p \in P\) (no. 2, Example 2) and, for p = 0, \(M_{p}\) and \(M_{p}^{**}\) are both equal to V.

THEOREM 4. Let E be a finitely generated A-module, T the torsion submodule of E and M = E/T. There exists a pseudo-isomorphism

\[
f \colon \mathbf {E} \rightarrow \mathbf {T} \times \mathbf {M}.
\]

We shall first show two lemmas.

LEMMA 2. Let \((\mathfrak{p}_i)_{1 \leqslant i \leqslant k}\) be a non-empty finite family of prime ideals of A of height 1 and let \(S = \bigcap_{i} (A - \mathfrak{p}_i)\) ; then the ring \(S^{-1}A\) is a principal ideal domain.

\(S^{-1}A\) is a semi-local ring whose maximal ideals are the \(m_{i} = p_{i}S^{-1}A\) for \(1 \leqslant i \leqslant k\) , the local ring \((S^{-1}A)_{m_{i}}\) being isomorphic to \(A_{p_{i}}\) (Chapter II, §3, no. 5, Proposition 17) and hence a discrete valuation ring. The ring \(S^{-1}A\) is therefore a Dedekind domain (§2, no. 2, Theorem 1 (f) and, as it is semi-local it is a principal ideal domain (§2, no. 2, Proposition 1).

LEMMA 3. There exists a homomorphism \(g: \mathbf{E} \to \mathbf{T}\) whose restriction to \(\mathbf{T}\) is both a homothety and a pseudo-isomorphism.

Let \(\mathfrak{a}\) be the annihilator of \(\mathrm{T}\) ; as \(\mathrm{T}\) is a finitely generated torsion A-module, \(\mathfrak{a} \neq 0\) . Let \(\mathfrak{p}_i\) ( \(1 \leqslant i \leqslant k\) ) be the prime ideals of height 1 containing \(\mathfrak{a}\) (which are finite in number (§ 1, no. 6, Theorem 4)); if this number is 0, T is pseudo-zero (Proposition 9(a)) and we may take g = 0. Otherwise, let \(S = \bigcap_{i} (A - p_{i})\) ; by Lemma 2, \(S^{-1}A\) is a principal ideal domain and hence \(S^{-1}M\) , which is a torsion-free finitely generated \(S^{-1}A\) -module, is free (Algebra, Chapter VII, § 4, no. 3, Corollary 2 to Theorem 2) and, as \(S^{-1}M = (S^{-1}E)/(S^{-1}T)\) , \(S^{-1}T\) is a direct factor of \(S^{-1}E\) (Algebra, Chapter II, § 1, no. 11, Proposition 21). Now,

\[
\operatorname{Hom} _ {\mathbf {S} ^ {- 1} \mathbf {A}} (\mathrm{S} ^ {- 1} \mathrm{E}, \mathrm{S} ^ {- 1} \mathrm{T}) = \mathrm{S} ^ {- 1} \operatorname{Hom} _ {\mathbf {A}} (\mathrm{E}, \mathrm{T})
\]

(Chapter II, §2, no. 7, Proposition 19); hence there exist \(s_0 \in S\) and \(g_0 \in \mathrm{Hom}_{\mathbf{A}}(\mathbf{E}, \mathbf{T})\) such that \(s_0^{-1}g_0\) is a projector of \(\mathbf{S}^{-1}\mathbf{E}\) onto \(\mathbf{S}^{-1}\mathbf{T}\) . If \(h_0 \in \mathrm{Hom}_{\mathbf{A}}(\mathbf{T}, \mathbf{T})\) denotes the restriction of \(g_0\) to \(\mathbf{T}\) , there therefore exists \(s_1 \in \mathbf{S}\) such that \(s_1h_0(x) = s_1s_0x\) for all \(x \in \mathbf{T}\) ; writing \(s = s_1s_0\) , \(g = s_1g_0\) , \(h = s_1h_0\) , \(h\) is therefore the homothety of ratio \(s\) on \(\mathbf{T}\) and is the restriction of \(g\) to \(\mathbf{T}\) . It remains to verify that \(h\) is a pseudo-isomorphism. Now, if \(p = 0\) or if \(p \in P\) is distinct from the \(p_i (1 \leqslant i \leqslant k)\) , \(T_p = 0\) (Chapter II, §4, no. 4, Proposition 17) and \(h_p: T_p \to T_p\) is an isomorphism; if on the contrary \(p\) is equal to one of the \(p_i (1 \leqslant i \leqslant k)\) , \(s\) is invertible in \(A_{p_i}\) and \(h_{p_i}\) , the homothety of ratio \(s\) on \(T_{p_i}\) , is also an isomorphism, which completes the proof of Lemma 3.

We now prove Theorem 4. Let \(g: \mathbf{E} \to \mathbf{T}\) be a homomorphism satisfying the properties of Lemma 3; let \(h\) be the restriction of \(g\) to \(\mathbf{T}\) and let \(\pi\) be the canonical projection of \(\mathbf{E}\) onto \(\mathbf{M}\) . We show that the homomorphism \(f = (g, \pi): \mathbf{E} \to \mathbf{T} \times \mathbf{M}\) solves the problem. There is the commutative diagram:

\[
\begin{array}{c c c} 0 \longrightarrow \mathrm{T} \longrightarrow & \mathrm{E} & \longrightarrow \mathrm{M} \longrightarrow 0 \\ h \Big \downarrow & f \Big \downarrow & ^ {1 _ {\mathrm{M}}} \Big \downarrow \\ 0 \longrightarrow \mathrm{T} \longrightarrow \mathrm{T} \times \mathrm{M} \longrightarrow \mathrm{M} \longrightarrow 0 \end{array}
\]

where the rows are exact. The snake diagram (Chapter I, § 1, no. 4, Proposition 2) gives the exact sequence:

\[
0 \rightarrow \operatorname{Ker} (h) \rightarrow \operatorname{Ker} (f) \rightarrow 0 \rightarrow \operatorname{Coker} (h) \rightarrow \operatorname{Coker} (f) \rightarrow 0
\]

and hence \(\operatorname{Ker}(f)\) is isomorphic to \(\operatorname{Ker}(h)\) and \(\operatorname{Coker}(f)\) to \(\operatorname{Coker}(h)\) . As \(h\) is a pseudo-isomorphism, so is \(f\) .

We can say that ``to within a pseudo-isomorphism'' Theorem 4 reduces the study of finitely generated A-modules to that of torsion-free modules on the one hand and to that of torsion modules on the other. Moreover, we have seen above (Example 3) that a torsion-free module is pseudo-isomorphic to its bidual and hence to a reflexive module. As for torsion modules, there is the following result, which determines them to within a pseudo-isomorphism:

THEOREM 5. Let T be a finitely generated torsion A-module. There exist two finite families \((n_i)_{i\in I}\) and \((p_i)_{i\in I}\) , where the \(n_i\) are integers \(\geqslant 1\) and the \(p_i\) are prime ideals of A of height 1 such that, if we write \(\mathrm{T}' = \bigoplus_{i\in I} \mathrm{A} / p_i^{n_i}\) , there exists a pseudo-isomorphism of T to \(\mathrm{T}'\) . Moreover, the families \((n_i)_{i\in I}\) and \((p_i)_{i\in I}\) with this property are unique to within a bijection of the indexing set and the \(p_i\) contain the annihilator of T.

Uniqueness: If \(f\colon T\to T'\) is a pseudo-isomorphism and \(\mathfrak{p}\in\mathbf{P},f_{\mathfrak{p}}\colon T_{\mathfrak{p}}\to T'_{\mathfrak{p}}\) is an isomorphism. Now, \(T'_{\mathfrak{p}}\) is the direct sum of the \(A_{\mathfrak{p}}/p^{n_{i}}A_{\mathfrak{p}}\) , the sum being over the indices \(i\) such that \(\mathfrak{p}_{i}=\mathfrak{p}\) ; the \(p^{n_{i}}A_{\mathfrak{p}}\) are therefore the elementary divisors of the torsion \(A_{\mathfrak{p}}\) -module \(T_{\mathfrak{p}}\) (Algebra, Chapter VII, § 4, no. 7); their uniqueness has been proved in Algebra, Chapter VII, § 4, no. 7, Proposition 7.

Existence: We may confine our attention to the case where \(\mathbf{T} \neq 0\) . Let \(\mathfrak{a}\) be the annihilator (non-zero and distinct from A) of \(\mathbf{T}\) , \(\mathfrak{p}_i\) ( \(1 \leqslant i \leqslant k\) ) the prime ideals of A of height 1 containing \(\mathfrak{a}\) (which are finite in number (§ 1, no. 6, Theorem 4)) and \(\mathbf{S} = \bigcap_i (\mathbf{A} - \mathfrak{p}_i)\) . The semi-local ring \(\mathbf{A}' = \mathbf{S}^{-1}\mathbf{A}\) is a principal ideal domain (Lemma 2) and has maximal ideals the \(\mathfrak{m}_i = \mathfrak{p}_i\mathbf{A}'\) ; as \(\mathbf{S}^{-1}\mathbf{T}\) is a finitely generated torsion \(\mathbf{A}'\) -module, it is isomorphic to a finite direct sum \(\bigoplus_{f \in \mathbf{I}} \mathbf{A}' / \mathfrak{m}_{\phi(f)}^{n_f}\) , where \(\phi\) is a mapping of a finite set I to \([1, k]\) (Algebra, Chapter VII, § 4, no. 7, Proposition 7); as \(\mathbf{A}' / \mathfrak{m}_{\phi(f)}^{n_f}\) is isomorphic to \(\mathbf{S}^{-1}(\mathbf{A} / \mathfrak{p}_{\phi(f)}^{n_f})\) (Chapter II, § 2, no. 4), we have obtained a torsion A-module \(\mathbf{T}'\) of the desired type and an isomorphism \(f_0\) of \(\mathbf{S}^{-1}\mathbf{T}\) onto \(\mathbf{S}^{-1}\mathbf{T}'\) . As \(\operatorname{Hom}_{\mathbf{S}^{-1}\mathbf{A}}(\mathbf{S}^{-1}\mathbf{T}, \mathbf{S}^{-1}\mathbf{T}')\) is equal to \(\mathbf{S}^{-1}\operatorname{Hom}_{\mathbf{A}}(\mathbf{T}, \mathbf{T}')\) (Chapter II, § 2, no. 7, Proposition 19), there exist \(s \in S\) and a homomorphism \(f: T \to T'\) such that \(f_0 = s^{-1}f\) . It remains to show that \(f\) is a pseudo-isomorphism: now, if \(\mathfrak{p} = 0\) or if \(\mathfrak{p} \in P\) is distinct from the \(\mathfrak{p}_i\) , \(\mathrm{T}_{\mathfrak{p}} = \mathrm{T}_{\mathfrak{p}}' = 0\) (Chapter II, § 4, no. 4, Proposition 17); if on the contrary \(\mathfrak{p}\) is one of the \(\mathfrak{p}_i\) ( \(1 \leqslant i \leqslant k\) ), \(s\) is invertible in \(\mathbf{A}_{\mathfrak{p}_i}\) and, as \(f_{\mathfrak{p}_i} = s(f_0)_{\mathfrak{p}_i}\) and \((f_0)_{\mathfrak{p}_i}\) is an isomorphism of \(\mathrm{T}_{\mathfrak{p}_i} = (\mathrm{S}^{-1}\mathrm{T})_{\mathfrak{m}_i}\) onto \(\mathrm{T}_{\mathfrak{p}_i}' = (\mathrm{S}^{-1}\mathrm{T}')_{\mathfrak{m}_i}\) , so is \(f_{\mathfrak{p}_i}\) .

Remark (2) In the statement of Theorem 5, the modules \(\mathbf{A} / \mathfrak{p}_i^n\) may be replaced by \(\mathbf{A} / \mathfrak{p}_i^{(n_i)}\) (§ 1, no. 4, Proposition 8). For all \(\mathfrak{p} \in \mathbb{P}\) , the canonical mapping \(g\colon \mathbf{A} / \mathfrak{p}^n \to \mathbf{A} / \mathfrak{p}^{(n)} = \mathbf{A} / (\mathbf{A} \cap \mathfrak{p}^n \mathbf{A}_{\mathfrak{p}})\) is a pseudo-isomorphism, as, for all \(\mathfrak{q} \in \mathbb{P}\) distinct from \(\mathfrak{p}\) , \(\mathbf{A}_{\mathfrak{q}} / \mathfrak{p}^n \mathbf{A}_{\mathfrak{q}} = \mathbf{A}_{\mathfrak{q}} / \mathfrak{p}^{(n)} \mathbf{A}_{\mathfrak{q}} = 0\) and \(\mathbf{A}_{\mathfrak{p}} / \mathfrak{p}^n \mathbf{A}_{\mathfrak{p}} = \mathbf{A}_{\mathfrak{p}} / \mathfrak{p}^{(n)} \mathbf{A}_{\mathfrak{p}}\) .

* Given an exact sequence of A-modules, E → F → G, if E and G are pseudo-zero, so is F, as follows from Definition 2 and Chapter II, § 2, no. 4, Theorem 1. In the language of categories, we may then say that, in the category C of A-modules, the sub-category C' of pseudo-zero modules is full and we may then define the quotient category C/C': the objects in this category are also A-modules but the set of morphisms from E to F (for E, F in C/C') is the direct limit of the set of commutative groups \(\mathrm{Hom}_{\mathbf{A}}(\mathbf{E}',\mathbf{F}')\) , where \(\mathbf{E}'\) (resp. \(\mathbf{F}'\) ) runs through the set of submodules of \(\mathbf{E}\) (resp. the set of quotient modules \(\mathbf{F} / \mathbf{F}''\) of \(\mathbf{F}\) ) such that \(\mathbf{E} / \mathbf{E}'\) (resp. \(\mathbf{F}''\) ) is pseudo-zero. Of course, for every ordered pair of A-modules \(\mathbf{E}\) , \(\mathbf{F}\) , there is a canonical homomorphism \(\mathrm{Hom}_{\mathcal{C}}(\mathbf{E},\mathbf{F})\to \mathrm{Hom}_{\mathcal{C} / \mathcal{C}'}(\mathbf{E},\mathbf{F})\) . To say that a homomorphism \(u\in \mathrm{Hom}_{\mathbf{A}}(\mathbf{E},\mathbf{F})\) is pseudo-zero (resp. pseudo-injective, pseudo-surjective, pseudo-bijective) means that its canonical image in \(\mathrm{Hom}_{\mathcal{C} / \mathcal{C}'}(\mathbf{E},\mathbf{F})\) is zero (resp. a monomorphism, an epimorphism, an isomorphism).*

\subsection{Divisors attached to torsion modules}

We keep the same notation and hypotheses as in nos. 2, 3 and 4. Recall that D(A) (or simply D) denotes the divisor group of A, written additively: we know (§ 1, no. 3, Theorem 2) that D is the free Z-module generated by the elements of P.

Let T be a finitely generated torsion A-module. For all \(p \in P\) , \(T_{p}\) is a finitely generated torsion \(A_{p}\) -module and hence a module of finite length (Chapter IV, §2, no. 5, Corollary 2 to Proposition 7); we shall denote this length by \(l_{p}(T)\) . Now \(T_{p} = 0\) for all p not containing the annihilator of T and hence for almost all p (§1, no. 6, Theorem 4), which justifies the following definition:

DEFINITION 4. If T is a finitely generated torsion A-module, the divisor:

\[
\chi (\mathbf {T}) = \sum_ {\mathfrak {p} \in \mathbf {P}} l _ {\mathfrak {p}} (\mathbf {T}). \mathfrak {p}.
\]

is called the content of T.

PROPOSITION 10. (i) Let \(0 \to T_1 \to T_2 \to T_3 \to 0\) be an exact sequence of finitely generated torsion A-modules. Then

\[
\chi (\mathrm{T} _ {2}) = \chi (\mathrm{T} _ {1}) + \chi (\mathrm{T} _ {3}).
\]

\begin{enumerate}
\def\labelenumi{(\roman{enumi})}
\setcounter{enumi}{1}
\item
  If there exists a pseudo-isomorphism \(f \colon \mathrm{T}_1 \to \mathrm{T}_2\) , then \(\chi(\mathrm{T}_1) = \chi(\mathrm{T}_2)\) .
\item
  In order that \(\chi(T) = 0\) , it is necessary and sufficient that \(T\) be pseudo-zero.
\end{enumerate}

In view of Definition 4, it suffices to consider for each \(p \in P\) the values of \(l_{p}\) for the torsion modules considered. Property (i) then follows from Chapter II, §2, no. 4, Theorem 1 and the additivity of lengths in an exact sequence (Algebra Chapter II, §1, no. 10, Proposition 16) and properties (ii) and (iii) follow immediately from the definitions in no. 4.

COROLLARY. Let \(0 \to T_n \to T_{n-1} \to \cdots \to T_0 \to 0\) be an exact sequence of finitely generated torsion A-modules. Then \(\sum_{i=0}^{n} (-1)^i \chi(T_i) = 0\) .

In view of Chapter II, §2, no. 4, Theorem 1, this follows again from the analogous property of the \(l_{\mathfrak{p}}\) (Algebra, Chapter II, §1, no. 10, Corollary 3 to Proposition 16).

Recall (Chapter II, § 5, no. 4) that we may speak of the set \(\mathbf{F}(\mathbf{A})\) of classes of finitely generated A-modules with respect to the relation of isomorphism; for every finitely generated A-module M, let \(\operatorname{cl}(\mathbf{M})\) denote the corresponding element of \(\mathbf{F}(\mathbf{A})\) ; we shall denote by \(\mathbf{T}(\mathbf{A})\) the subset of \(\mathbf{F}(\mathbf{A})\) consisting of the classes of finitely generated torsion A-modules. Clearly \(\chi\) defines a mapping of \(\mathbf{T}(\mathbf{A})\) to \(\mathbf{D}(\mathbf{A})\) , also denoted by \(\chi\) , such that \(\chi(\operatorname{cl}(\mathbf{T})) = \chi(\mathbf{T})\) .

PROPOSITION 11. Let G be a commutative group written additively and \(\phi: T(A) \to G\) a mapping; for every finitely generated torsion A-module T, we also write, by an abuse of language, \(\phi(T) = \phi(\text{cl}(T))\) . Suppose that the following conditions are satisfied:

\begin{enumerate}
\def\labelenumi{(\arabic{enumi})}
\item
  If \(0 \to \mathrm{T}_1 \to \mathrm{T}_2 \to \mathrm{T}_3 \to 0\) is an exact sequence of finitely generated torsion A-modules, then \(\phi(\mathrm{T}_2) = \phi(\mathrm{T}_1) + \phi(\mathrm{T}_3)\) .
\item
  If \(\mathbf{T}\) is pseudo-zero, then \(\phi (\mathbf{T}) = 0\) .
\end{enumerate}

Then there exists a unique homomorphism \(\theta: \mathbf{D}(\mathbf{A}) \to \mathbf{G}\) such that \(\phi = \theta \circ \chi\) .

As \(\chi(A/p) = p\) for all \(p\) , necessarily \(\theta(p) = \phi(A/p)\) for all \(p \in P\) , which proves the uniqueness of \(\theta\) , since the elements of \(P\) form a basis of \(D(A)\) . Conversely, let \(\theta\) be the homomorphism from \(D(A)\) to \(G\) such that \(\theta(p) = \phi(A/p)\) for all \(p \in P\) and let us show that it solves the problem. For this, we write

\[
\psi (\mathbf {T}) = \phi (\mathbf {T}) - \theta (\chi (\mathbf {T}))
\]

for every finitely generated torsion A-module T; clearly conditions (1) and (2) are also satisfied if \(\phi\) is replaced by \(\psi\) . On the other hand, \(\psi(A/p) = 0\) if \(p \in P\) ; if \(p\) is a prime ideal \(\neq 0\) and not in \(P\) , the annihilator of \(A/p\) is contained in no ideal of \(P\) , hence (no. 4, Theorem 5) \(A/p\) is pseudo-zero and therefore \(\psi(A/p) = 0\) . This being so, every finitely generated torsion A-module T admits a decomposition series whose factors are isomorphic to A-modules of the form \(A/p\) , where \(p \in \text{Supp}(T)\) (Chapter IV, § 1, no. 4, Theorems 1 and 2), and hence \(p \neq 0\) since \(T\) is a torsion module. By induction on the length of this decomposition series, we deduce (in view of property (1) for \(\psi\) ) that \(\psi(T) = 0\) .

* We may, as in no. 4, consider the quotient category \(\mathcal{T} / \mathcal{T}'\) of the category \(\mathcal{T}\) of finitely generated torsion A-modules by the full sub-category \(\mathcal{T}'\) of pseudo-zero finitely generated torsion A-modules. In the language of Abelian categories, Proposition 11 then expresses the fact that the Grothendieck group of the Abelian category \(\mathcal{T} / \mathcal{T}'\) is canonically isomorphic to D(A). *

PROPOSITION 12. If a is an ideal \(\neq 0\) of A,

\[
\chi (\mathrm{A} / \mathfrak {a}) = \chi ((\mathrm{A}: \mathfrak {a}) / \mathrm{A}) = \operatorname{div} \mathfrak {a}.
\]

Let \(p \in P\) . Then \(aA_p = p^{n_p}A_p\) where \(n_p \geqslant 0\) , since \(A_p\) is a discrete valuation ring. As \((A/a)_p = A_p/aA_p\) , \(l_p(A/a) = n_p\) , whence \(\chi(A/a) = \sum_{p \in P} n_p p = \text{div } a\) ( \(\S 1, \text{ no. } 4, \text{ Proposition } 7\) ).

On the other hand, \((\mathbf{A}:\mathfrak{a})_{\mathfrak{p}} = \mathbf{A}_{\mathfrak{p}}\colon \mathfrak{a}\mathbf{A}_{\mathfrak{p}} = \mathfrak{p}^{-n_{\mathfrak{p}}}\mathbf{A}_{\mathfrak{p}}\) , hence \(l_{\mathfrak{p}}((\mathbf{A}:\mathfrak{a}) / \mathbf{A}) = n_{\mathfrak{p}}\) and we conclude in the same way.

\subsection{Relative invariant of two lattices}

We keep the notation and hypotheses of nos. 2 to 5. Let V be a vector space of finite rank n over K and M a lattice of V with respect to A. Let W be the exterior power \(\bigwedge^{n}V\) , which is a vector space of rank 1 over K and let \(M_{W}\) denote the lattice of W generated by the image of \(M^{n}\) under the canonical mapping \(V^{n}\to\bigwedge^{n}V\) (no. 1, Proposition 3 (iii); note that \(M_{W}\) is not necessarily isomorphic to \(\bigwedge^{n}M\) (Algebra, Chapter III, § 5, Exercise 9)). If e is a basis of W over K, we may write \(M_{W}=a.e\) , where a is a fractional ideal \(\neq0\) of A.

Let \(\mathbf{M}'\) be another lattice of \(\mathbf{V}\) and let us write \(\mathbf{M}_{\mathbf{W}}' = \mathfrak{a}'\) .e, where \(\mathfrak{a}'\) is a fractional ideal \(\neq 0\) of \(\mathbf{A}\) ; the divisor \(\operatorname{div}(\mathfrak{a}) - \operatorname{div}(\mathfrak{a}')\) does not depend on the choice of basis \(e\) of \(\mathbf{W}\) , \(\mathfrak{a}\) and \(\mathfrak{a}'\) being multiplied by the same element of \(\mathbf{K}^*\) when the basis is changed; we shall write \(\chi(\mathbf{M}, \mathbf{M}') = \operatorname{div}(\mathfrak{a}) - \operatorname{div}(\mathfrak{a}')\) and say that this divisor is the relative invariant of \(\mathbf{M}'\) with respect to \(\mathbf{M}\) . Clearly, if \(\mathbf{M}, \mathbf{M}', \mathbf{M}''\) are three lattices of \(\mathbf{V}\) , then:

\[
\chi (\mathbf {M}, \mathbf {M} ^ {\prime}) + \chi (\mathbf {M} ^ {\prime}, \mathbf {M} ^ {\prime \prime}) + \chi (\mathbf {M} ^ {\prime \prime}, \mathbf {M}) = 0\tag{3}
\]

\[
\chi (\mathbf {M}, \mathbf {M} ^ {\prime}) + \chi (\mathbf {M} ^ {\prime}, \mathbf {M}) = 0\tag{4}.
\]

For all \(p \in P\) , it follows immediately from the definitions that \((M_{W})_p = (M_p)_{W}\) ; moreover, since \(M_p\) is then a free \(A_p\) -module since \(A_p\) is a principal ideal domain, a basis of \(M_p\) over \(A_p\) is a basis of \(V\) over \(K\) , hence \((M_p)_W = \bigwedge^n (M_p)\) (Chapter II, § 2, no. 8) and the fractional ideal \(a_p = aA_p\) is a principal ideal domain. If we set \(a_p = p^{n_p} A_p\) , \(a' = p^{n'_{p}} A_p\) , then:

\[
\chi (\mathbf {M}, \mathbf {M} ^ {\prime}) = \sum_ {p \in P} (n _ {p} - n _ {p} ^ {\prime}). p,
\]

which may also be written as:

\[
\chi (\mathbf {M}, \mathbf {M} ^ {\prime}) = \sum_ {p \in P} \chi (\mathbf {M} _ {p}, \mathbf {M} _ {p} ^ {\prime})\tag{5}
\]

identifying \(\mathbf{D}(\mathbf{A}_{\mathfrak{p}})\) with the sub-Z-module of \(\mathbf{D}(\mathbf{A})\) generated by p.

PROPOSITION 13. Let M be a lattice of V and u a K-automorphism of V. Then:

\[
- \chi (\mathbf {M}, u (\mathbf {M})) = \operatorname{div} (\det (u)).\tag{6}
\]

For all \(p \in P\) , then \(\bigwedge^{n} (u(M)_{p}) = \bigwedge^{n} (u(M_{p}))\) ; if \((e_{i})_{1 \leqslant i \leqslant n}\) is a basis of \(M_{p}\) , then

\[
\bigwedge^ {n} \left(\mathrm{M} _ {p}\right) = \mathrm{A} _ {p}. e _ {1} \wedge e _ {2} \wedge \dots \wedge e _ {n},
\]

and \(\bigwedge^{n}(u(\mathbf{M}_{\mathfrak{p}})) = \mathbf{A}_{\mathfrak{p}}.\det (u)e_{1}\wedge e_{2}\wedge \dots \wedge e_{n}\) , whence the proposition by virtue of formula (5).

PROPOSITION 14. If M, M' are two lattices of V such that M' ⊂ M, M/M' is a finitely generated torsion A-module and:

\[
\chi (\mathbf {M}, \mathbf {M} ^ {\prime}) = - \chi (\mathbf {M} / \mathbf {M} ^ {\prime}).\tag{7}
\]

Clearly \(\mathbf{M} / \mathbf{M}' \subset \mathbf{V} / \mathbf{M}'\) is a finitely generated torsion module; on the other hand, for all \(\mathfrak{p} \in \mathbb{P}\) , we know (Algebra, Chapter VII, § 4, no. 2, Theorem 1) that there exist bases \((e_i)_{1 \leqslant i \leqslant n}\) of \(\mathbf{M}_{\mathfrak{p}}\) and \((e_i')_{1 \leqslant i \leqslant n}\) of \(\mathbf{M}_{\mathfrak{p}}'\) such that \(e_i' = \pi^{\nu_i} e_i\) for \(1 \leqslant i \leqslant n\) and integers \(\nu_i \geqslant 0\) , \(\pi\) being a uniformizer of \(\mathbf{A}_{\mathfrak{p}}\) . Therefore (in the notation introduced above) \(n_{\mathfrak{p}}' - n_{\mathfrak{p}} = \sum_{i=1}^{n} \nu_i\) ; and on the other hand, \((\mathbf{M} / \mathbf{M}')_{\mathfrak{p}} = \mathbf{M}_{\mathfrak{p}} / \mathbf{M}_{\mathfrak{p}}'\) is isomorphic to the torsion \(\mathbf{A}_{\mathfrak{p}}\) -module \(\bigoplus_{i=1}^{n} \mathbf{A}_{\mathfrak{p}} / p^{\nu_i} \mathbf{A}_{\mathfrak{p}}\) and hence its length is \(\sum_{i=1}^{n} \nu_i\) , which proves the proposition, in view of (5) and Definition 4 of no. 5.

COROLLARY. Let \(L_{1}\) , \(L_{2}\) be two free A-modules of the same rank n and let \(f: L_{1} \to L_{2}\) be a homomorphism. Let U be the matrix of f with respect to bases of \(L_{1}\) and \(L_{2}\) . For Coker(f) to be a torsion A-module, it is necessary and sufficient that \(\det(U) \neq 0\) and then:

\[
\chi (\operatorname{Coker} (f)) = \operatorname{div} (\det (U)).\tag{8}
\]

\(L_{1}\) and \(L_{2}\) can be considered as lattices in \(V_{1} = L_{1} \otimes_{A} K\) and \(V_{2} = L_{2} \otimes_{A} K\) respectively, f extending to a K-homomorphism \(f_{(K)}: V_{1} \to V_{2}\) . Then

\[
\left(\operatorname{Coker} (f)\right) _ {(\mathbf {K})} = \operatorname{Coker} \left(f _ {(\mathbf {K})}\right)
\]

and to say that \(\operatorname{Coker}(f)\) is a torsion A-module means that \(\operatorname{Coker}(f_{(\mathbf{K})}) = 0\) ; now, it amounts to the same to say that \(f_{(\mathbf{K})}\) is surjective or that \(\det(U) \neq 0\) , whence the first assertion. On the other hand, we may write \(f(\mathbf{L}_1) = u(\mathbf{L}_2)\) , where \(u\) is an endomorphism of \(\mathbf{L}_2\) of determinant \(\det(U)\) ; as

\[
\operatorname{Coker} (f) = \mathrm{L} _ {2} / u (\mathrm{L} _ {2}),
\]

formula (8) follows from (7) and (6).

Example. If A = Z, the divisor group of A is identified with the multiplicative group \(Q_{+}^{*}\) of rational numbers \textgreater0. For every finite commutative group T, \(\chi(T)\) is the order of T; the above corollary shows that the order of the group \(\operatorname{Coker}(f)\) is equal to the absolute value of \(\det(U)\) (cf.~Algebra, Chapter VII, §4, no. 7, Corollary 3 to Theorem 3).
